\documentclass[9pt,a4paper,landscape]{article}
\usepackage[margin=1.4cm,landscape]{geometry}
\usepackage{fontspec,longtable,array,booktabs,xcolor}
\setmainfont{Latin Modern Roman}
\newfontfamily\hifont{Nirmala UI}
\definecolor{hdr}{HTML}{1F3864}
\setlength{\tabcolsep}{4pt}\renewcommand{\arraystretch}{1.18}
\usepackage{fancyhdr}\pagestyle{fancy}\fancyhf{}
\fancyhead[L]{\small Kedarnath Yatra Worker Survey --- Question Register}
\fancyhead[R]{\small\thepage}\renewcommand{\headrulewidth}{0.3pt}
\begin{document}
\begin{center}
{\LARGE\bfseries Kedarnath Yatra Worker Survey}\\[3pt]
{\large Question register, with sources}\\[8pt]
193 questions asked \quad$\cdot$\quad 143 variables constructed \quad$\cdot$\quad 148 distinct sources\\[3pt]
{\small Every question is shown in English and Hindi exactly as the tablet presents it. The source
column names the instrument or paper the item is taken from or adapted from.}
\end{center}
\vspace{4pt}
\noindent{\large\bfseries\color{hdr} Module A --- Respondent and household}\quad{\small(20 questions)}\\[3pt]
{\small\itshape Read aloud before this module: First, a few things about you and the people who live in your household.}\\[5pt]
\begin{longtable}{@{}p{0.5cm}p{3.3cm}p{1.7cm}p{7.3cm}p{5.6cm}p{5.4cm}@{}}
\toprule \footnotesize\textbf{\#} & \footnotesize\textbf{Variable} & \footnotesize\textbf{Type} & \footnotesize\textbf{Question (English / Hindi)} & \footnotesize\textbf{Options {\&} skip} & \footnotesize\textbf{Source}\\ \midrule
\endhead
\footnotesize 1 & \footnotesize\texttt{age} & \scriptsize number & \footnotesize How old are you? (completed years)\par\vspace{1pt}{\footnotesize\color{hdr}{\hifont आपकी उम्र कितनी है? (पूरे साल)}} & \scriptsize  & \scriptsize Standard; sensitivity covariate in the VEP model (Azeem et al. 2016 structure)\\
\footnotesize 2 & \footnotesize\texttt{female} & \scriptsize yes/no & \footnotesize Record sex; ask if unsure: are you male or female?\par\vspace{1pt}{\footnotesize\color{hdr}{\hifont आप पुरुष हैं या महिला?}} & \scriptsize 0 Male | 1 Female & \scriptsize Standard household-survey item\\
\footnotesize 3 & \footnotesize\texttt{hoh\_relation} & \scriptsize select one & \footnotesize Who is the head of your household — is it you, or someone else? If someone else, how are they related to you?\par\vspace{1pt}{\footnotesize\color{hdr}{\hifont आपके घर का मुखिया कौन है — आप खुद, या कोई और? अगर कोई और है, तो वे आपके क्या लगते हैं?}} & \scriptsize 1 Self (respondent is the head) | 2 Husband | 3 Wife | 4 Father | 5 Mother | 6 Son | 7 Daughter | 8 Brother | 9 Sister | 10 Other male relative | 11 Other female relative | 12 Other, non-relative (male) | 13 Other, non-relative (female) & \scriptsize Standard household-survey item\\
\footnotesize 4 & \footnotesize\texttt{native\_language} & \scriptsize select one & \footnotesize What is your native language, the language you first learned at home?\par\vspace{1pt}{\footnotesize\color{hdr}{\hifont आपकी मातृभाषा कौन सी है — वह भाषा जो आपने घर में सबसे पहले सीखी?}} & \scriptsize 1 Garhwali | 2 Kumaoni | 3 Hindi | 4 Nepali | 5 Bhojpuri | 6 Assamese | 7 Bengali | 8 Bodo | 9 Dogri | 10 Gujarati | 11 Kannada | 12 Kashmiri | 13 Konkani | 14 Maithili | 15 Malayalam | 16 Manipuri (Meitei) | 17 Marathi | 18 Odia | 19 Punjabi | 20 Sanskrit | 21 Santali | 22 Sindhi | 23 Tamil | 24 Telugu | 25 Urdu | 96 Other Indian language | 97 Other / foreign language & \scriptsize Project design (language proxy for regional/socioeconomic origin; avoids a sensitive caste question and covers out-of-country migrants)\\
\footnotesize 5 & \footnotesize\texttt{native\_language\_other} & \scriptsize free text & \footnotesize Which language is it? Write down exactly what they say.\par\vspace{1pt}{\footnotesize\color{hdr}{\hifont वह कौन सी भाषा है? जैसा वे कहें वैसा ही लिख लीजिए।}} & \scriptsize \textit{Asked only if: Ask only if native\_language is 96 (other Indian language) or 97 (other / foreign language)} & \scriptsize Project design (the language item stands in for caste; an uncoded 'other' is a real loss)\\
\footnotesize 6 & \footnotesize\texttt{origin} & \scriptsize select one & \footnotesize Where is your permanent home?\par\vspace{1pt}{\footnotesize\color{hdr}{\hifont आपका पक्का घर कहाँ है?}} & \scriptsize 1 Local (same district) | 2 Other Uttarakhand district | 3 Other Indian state | 4 Nepal & \scriptsize Project design (distance bucket); NOT an NSS classification — the earlier citation to one was withdrawn 2026-09-28, no NSS schedule is held by this project\\
\footnotesize 7 & \footnotesize\texttt{home\_state} & \scriptsize select one & \footnotesize Which state or union territory is your permanent home in?\par\vspace{1pt}{\footnotesize\color{hdr}{\hifont वह घर किस राज्य में है? (अगर भारत से बाहर है तो वह बताइए।)}} & \scriptsize 1 Andhra Pradesh | 2 Arunachal Pradesh | 3 Assam | 4 Bihar | 5 Chhattisgarh | 6 Goa | 7 Gujarat | 8 Haryana | 9 Himachal Pradesh | 10 Jharkhand | 11 Karnataka | 12 Kerala | 13 Madhya Pradesh | 14 Maharashtra | 15 Manipur | 16 Meghalaya | 17 Mizoram | 18 Nagaland | 19 Odisha | 20 Punjab | 21 Rajasthan | 22 Sikkim | 23 Tamil Nadu | 24 Telangana | 25 Tripura | 26 Uttar Pradesh | 28 West Bengal | 29 Andaman and Nicobar Islands | 30 Chandigarh | 31 Dadra and Nagar Haveli and Daman and Diu | 32 Delhi | 33 Jammu and Kashmir | 34 Ladakh | 35 Lakshadweep | 36 Puducherry & \scriptsize Standard state/UT classification; needed to apply the right poverty line per respondent\\
\footnotesize 8 & \footnotesize\texttt{home\_rural\_urban} & \scriptsize select one & \footnotesize Is that home in a village, or in a town or city?\par\vspace{1pt}{\footnotesize\color{hdr}{\hifont वह घर गाँव में है, या क़स्बे या शहर में?}} & \scriptsize 1 Village (rural) | 2 Town or city (urban) & \scriptsize Rural and urban poverty lines (Sethu et al. 2024)\\
\footnotesize 9 & \footnotesize\texttt{marital\_status} & \scriptsize select one & \footnotesize Are you currently married, never married, or widowed/divorced/separated?\par\vspace{1pt}{\footnotesize\color{hdr}{\hifont क्या आपकी शादी हुई है, नहीं हुई है, या आप विधवा/विधुर या अलग हैं?}} & \scriptsize 1 Currently married | 2 Never married | 3 Widowed/divorced/separated & \scriptsize Standard household-survey item\\
\footnotesize 10 & \footnotesize\texttt{knows\_years\_schooling} & \scriptsize yes/no & \footnotesize Do you know exactly how many years of schooling you completed, counting from class 1?\par\vspace{1pt}{\footnotesize\color{hdr}{\hifont क्या आपको ठीक-ठीक याद है कि आपने कितने साल पढ़ाई की, कक्षा 1 से गिनकर?}} & \scriptsize 0 No | 1 Yes & \scriptsize Standard household-survey item\\
\footnotesize 11 & \footnotesize\texttt{years\_schooling} & \scriptsize whole number & \footnotesize How many years of schooling did you complete in total? (0 if none)\par\vspace{1pt}{\footnotesize\color{hdr}{\hifont आपने कुल कितने साल पढ़ाई की? (अगर बिल्कुल नहीं पढ़े तो 0)}} & \scriptsize \textit{Asked only if: Ask only if knows\_years\_schooling = 1} & \scriptsize NITI Aayog National MPI (years of schooling); Chaudhuri et al. 2002 covariate\\
\footnotesize 12 & \footnotesize\texttt{education\_level\_cat} & \scriptsize select one & \footnotesize Which is closest: no formal education; some schooling but did not finish primary; finished primary but not secondary; or finished secondary or higher?\par\vspace{1pt}{\footnotesize\color{hdr}{\hifont इनमें से कौन सा सबसे नज़दीक है: बिल्कुल नहीं पढ़े; थोड़ा पढ़े पर प्राइमरी पूरी नहीं की; प्राइमरी पूरी की पर दसवीं नहीं; या दसवीं या उससे ज़्यादा?}} & \scriptsize 1 No formal education | 2 Some schooling, did not complete primary | 3 Completed primary but not secondary | 4 Completed secondary or higher | 98 Prefer not to say\par\vspace{1pt}\textit{Asked only if: Ask only if knows\_years\_schooling = 0} & \scriptsize NITI Aayog National MPI (years of schooling); Chaudhuri et al. 2002 covariate\\
\footnotesize 13 & \footnotesize\texttt{training\_received} & \scriptsize yes/no & \footnotesize Apart from school and college, have you ever completed any training course or apprenticeship?\par\vspace{1pt}{\footnotesize\color{hdr}{\hifont स्कूल और कॉलेज की पढ़ाई छोड़कर, क्या आपने कभी कोई ट्रेनिंग या सीखने का कोर्स पूरा किया है?}} & \scriptsize 0 No | 1 Yes & \scriptsize STEP module 2; Chaudhuri et al. 2002 adaptive-capacity covariate\\
\footnotesize 14 & \footnotesize\texttt{hhsize} & \scriptsize whole number & \footnotesize How many members live in your household, including you?\par\vspace{1pt}{\footnotesize\color{hdr}{\hifont आपके घर में आपको मिलाकर कुल कितने लोग रहते हैं?}} & \scriptsize  & \scriptsize Chaudhuri et al. 2002; used to get per-capita consumption\\
\footnotesize 15 & \footnotesize\texttt{any\_member\_6yr\_schooling} & \scriptsize yes/no & \footnotesize Counting everyone in your household aged 10 or above, has at least one of them finished six years of schooling or more?\par\vspace{1pt}{\footnotesize\color{hdr}{\hifont आपके घर में 10 साल या उससे बड़े सभी लोगों में से, क्या कम से कम एक ने छह साल या उससे ज़्यादा पढ़ाई पूरी की है?}} & \scriptsize 0 No | 1 Yes\par\vspace{1pt}\textit{Asked only if: Ask only if the respondent's own schooling does not already settle it: years\_schooling < 6, or the bracket is not 'completed secondary or higher'} & \scriptsize NITI Aayog National MPI, Years of Schooling (weight 1/6)\\
\footnotesize 16 & \footnotesize\texttt{n\_children\_u6} & \scriptsize whole number & \footnotesize How many of them are children under 6 years old?\par\vspace{1pt}{\footnotesize\color{hdr}{\hifont इनमें 6 साल से छोटे बच्चे कितने हैं?}} & \scriptsize  & \scriptsize Adult-equivalent scales (Claro et al. 2010; Awuni et al. 2023); NITI Aayog National MPI nutrition age range\\
\footnotesize 17 & \footnotesize\texttt{n\_children\_6\_14} & \scriptsize whole number & \footnotesize And how many are aged 6 to 14?\par\vspace{1pt}{\footnotesize\color{hdr}{\hifont और 6 से 14 साल के कितने हैं?}} & \scriptsize  & \scriptsize NITI Aayog National MPI school-attendance indicator\\
\footnotesize 18 & \footnotesize\texttt{n\_children\_in\_school} & \scriptsize whole number & \footnotesize How many of these children go to school now?\par\vspace{1pt}{\footnotesize\color{hdr}{\hifont इनमें से कितने बच्चे अभी स्कूल जाते हैं?}} & \scriptsize \textit{Asked only if: Ask only if n\_children\_6\_14 > 0} & \scriptsize NITI Aayog National MPI school-attendance indicator\\
\footnotesize 19 & \footnotesize\texttt{n\_earners} & \scriptsize whole number & \footnotesize How many of them, including you, earn money?\par\vspace{1pt}{\footnotesize\color{hdr}{\hifont इनमें से कितने लोग, आपको मिलाकर, पैसा कमाते हैं?}} & \scriptsize  & \scriptsize Dependency measure used in VEP studies (Azeem et al. 2016)\\
\footnotesize 20 & \footnotesize\texttt{main\_income\_earner} & \scriptsize yes/no & \footnotesize Among those who earn money, are you the main income earner in your household?\par\vspace{1pt}{\footnotesize\color{hdr}{\hifont जो लोग कमाते हैं, उनमें सबसे ज़्यादा कमाने वाले आप ही हैं?}} & \scriptsize 0 No | 1 Yes\par\vspace{1pt}\textit{Asked only if: Ask only if n\_earners > 1 (if the respondent is the only earner, this is automatic)} & \scriptsize Apablaza et al. 2026, Appendix 2 Q4 (adapted to a single-respondent report: 'who is the main contributor' becomes 'are you the main contributor')\\
\bottomrule\end{longtable}\vspace{6pt}
\noindent{\large\bfseries\color{hdr} Module B --- Work and work history}\quad{\small(15 questions)}\\[3pt]
{\small\itshape Read aloud before this module: Now about the work you do here on the Yatra route, and the work you did before this.}\\[5pt]
\begin{longtable}{@{}p{0.5cm}p{3.3cm}p{1.7cm}p{7.3cm}p{5.6cm}p{5.4cm}@{}}
\toprule \footnotesize\textbf{\#} & \footnotesize\textbf{Variable} & \footnotesize\textbf{Type} & \footnotesize\textbf{Question (English / Hindi)} & \footnotesize\textbf{Options {\&} skip} & \footnotesize\textbf{Source}\\ \midrule
\endhead
\footnotesize 21 & \footnotesize\texttt{occupation} & \scriptsize select one & \footnotesize What is your main work during the Yatra season? Ask what they do and code the closest group; do not read the list aloud unless needed.\par\vspace{1pt}{\footnotesize\color{hdr}{\hifont यात्रा के मौसम में आपका मुख्य काम क्या है?}} & \scriptsize 1 Pony/mule worker (works the animals, does not own them) | 2 Pony/mule owner (owns the animals) | 3 Porter (carries on own back/shoulders - goods, luggage, or a person in a kandi) | 4 Palki / dandi bearer (carries a person on a palanquin, as part of a team) | 5 Dhaba, tea stall or food stall WORKER (prepared food or tea) | 6 Dhaba, tea stall or food stall OWNER (prepared food or tea) | 7 Shop WORKER (goods: prasad, puja items, clothes, general store) | 8 Shop OWNER (goods, not prepared food) | 9 Hotel / lodge WORKER | 10 Hotel / lodge OWNER | 11 Driver (any motor vehicle) | 12 Guide | 13 Wage labourer (construction, loading, odd jobs) | 14 Other & \scriptsize Project quota design, mapped to NCO-2015 families\\
\footnotesize 22 & \footnotesize\texttt{occupation\_detail} & \scriptsize free text & \footnotesize In your own words, what exactly is your main work here? Write down what they say — what they actually do, and who for. Do not tick a box for this one.\par\vspace{1pt}{\footnotesize\color{hdr}{\hifont अपने शब्दों में बताइए, यहाँ आपका मुख्य काम ठीक-ठीक क्या है? आप करते क्या हैं, और किसके लिए?}} & \scriptsize  & \scriptsize Project design (coarse quota group + verbatim detail, coded to NCO-2015 in the office)\\
\footnotesize 23 & \footnotesize\texttt{employment\_type} & \scriptsize select one & \footnotesize In this main work, are you: working on your own account without hired workers; running a business with hired workers; a regular monthly wage-earner; or a daily/casual wage-earner?\par\vspace{1pt}{\footnotesize\color{hdr}{\hifont इस काम में आप क्या हैं: अपना काम खुद करते हैं और कोई नौकर नहीं रखते; अपना काम है और दूसरों को काम पर रखते हैं; महीने की पक्की तनख़्वाह पाते हैं; या दिहाड़ी पर काम करते हैं?}} & \scriptsize 1 Own-account (self-employed, no hired workers) | 2 Employer/business owner (self-employed) | 3 Regular wage or salary | 4 Casual or daily wage labour | 5 Unpaid family worker & \scriptsize PLFS / NSS employment status classification [source not held by this project; citation unverified as of 2026-09-28]\\
\footnotesize 24 & \footnotesize\texttt{wage\_employer} & \scriptsize select one & \footnotesize Who employs you? A private company; a government or public body; a household (domestic work); a contractor or agent; or an individual employer?\par\vspace{1pt}{\footnotesize\color{hdr}{\hifont आपको काम पर कौन रखता है? कोई निजी कंपनी; सरकारी या सार्वजनिक विभाग; कोई घर (घरेलू काम); ठेकेदार या एजेंट; या कोई व्यक्ति जो सीधे पैसे देता है?}} & \scriptsize 1 Private company | 2 Government or public body | 3 A household (domestic work) | 4 A contractor or agent | 5 An individual employer | 97 Don't know\par\vspace{1pt}\textit{Asked only if: Ask only if employment\_type is 3 or 4 (regular or casual wage worker)} & \scriptsize Apablaza et al. 2026, Appendix 2 Q8 (shortened to wage workers); categories ours\\
\footnotesize 25 & \footnotesize\texttt{other\_activity\_types} & \scriptsize select multiple & \footnotesize Besides your main work, what other paid work do you do during the Yatra season? Check every one that applies; leave blank if there is none. Ask and code the closest group for each; do not read the list aloud unless needed.\par\vspace{1pt}{\footnotesize\color{hdr}{\hifont वे कौन-कौन से काम हैं? जो-जो लागू हो, सब बताइए।}} & \scriptsize 1 Pony/mule worker (works the animals, does not own them) | 2 Pony/mule owner (owns the animals) | 3 Porter (carries on own back/shoulders - goods, luggage, or a person in a kandi) | 4 Palki / dandi bearer (carries a person on a palanquin, as part of a team) | 5 Dhaba, tea stall or food stall WORKER (prepared food or tea) | 6 Dhaba, tea stall or food stall OWNER (prepared food or tea) | 7 Shop WORKER (goods: prasad, puja items, clothes, general store) | 8 Shop OWNER (goods, not prepared food) | 9 Hotel / lodge WORKER | 10 Hotel / lodge OWNER | 11 Driver (any motor vehicle) | 12 Guide | 13 Wage labourer (construction, loading, odd jobs) | 14 Other\par\vspace{1pt}\textit{Asked only if: May be left blank} & \scriptsize Project design (multiple job-holding, common in the Yatra economy)\\
\footnotesize 26 & \footnotesize\texttt{other\_activity\_income\_pm} & \scriptsize amount (Rs) & \footnotesize Taking all of that other work together, about how much do you usually earn from it in a month? (Rs; leave blank if there is no other work)\par\vspace{1pt}{\footnotesize\color{hdr}{\hifont वह सारा दूसरा काम मिलाकर, महीने में लगभग कितना कमा लेते हैं? (रुपये)}} & \scriptsize \textit{Asked only if: May be left blank} & \scriptsize Project design (multiple job-holding)\\
\footnotesize 27 & \footnotesize\texttt{owns\_shop\_stall} & \scriptsize yes/no & \footnotesize Does your household own a shop or stall that you use to earn?\par\vspace{1pt}{\footnotesize\color{hdr}{\hifont क्या आपके घर की अपनी दुकान या खोखा है जिससे आप कमाते हैं?}} & \scriptsize 0 No | 1 Yes & \scriptsize Chaudhuri et al. 2002 asset covariate\\
\footnotesize 28 & \footnotesize\texttt{employers\_in\_season} & \scriptsize whole number & \footnotesize During one Yatra season, how many different employers or businesses do you work in? (1 if you stay with the same one, or run the same one, all season)\par\vspace{1pt}{\footnotesize\color{hdr}{\hifont एक यात्रा सीज़न में आप कितने अलग-अलग मालिकों के पास या कितने अलग-अलग कामों में काम करते हैं? (अगर पूरे सीज़न एक ही जगह रहते हैं या अपना एक ही काम चलाते हैं, तो 1 लिखें)}} & \scriptsize  & \scriptsize Apablaza et al. 2026 Appendix 2 Q7 (concept: job stability within the engagement); asked as a count per this project's no-scale rule\\
\footnotesize 29 & \footnotesize\texttt{years\_current\_job} & \scriptsize number & \footnotesize How many years have you been in this work? Write it as a decimal if needed, for example 1.5 for one year and six months, or 0.5 for six months.\par\vspace{1pt}{\footnotesize\color{hdr}{\hifont अगर एक साल या ज़्यादा: आप इस काम में कितने साल से हैं? दशमलव में लिखें, जैसे डेढ़ साल के लिए 1.5।}} & \scriptsize  & \scriptsize Apablaza et al. 2026, Appendix 2 Q10; Table 3 (stability)\\
\footnotesize 30 & \footnotesize\texttt{hours\_day\_yatra} & \scriptsize whole number & \footnotesize In the Yatra season, on a normal working day, how many hours do you work in total, across all your work?\par\vspace{1pt}{\footnotesize\color{hdr}{\hifont यात्रा के मौसम में, आम दिन में आप कुल कितने घंटे काम करते हैं, सारे कामों को मिलाकर?}} & \scriptsize  & \scriptsize Apablaza et al. 2026 Q13 (decomposed into hours/day x days/week rather than asked as one weekly figure, for ease of recall)\\
\footnotesize 31 & \footnotesize\texttt{days\_week\_yatra} & \scriptsize whole number & \footnotesize How many days a week do you usually work in the season?\par\vspace{1pt}{\footnotesize\color{hdr}{\hifont सीज़न में आप हफ़्ते में कितने दिन काम करते हैं?}} & \scriptsize  & \scriptsize Apablaza et al. 2026 Q13 (decomposed into hours/day x days/week rather than asked as one weekly figure, for ease of recall)\\
\footnotesize 32 & \footnotesize\texttt{prev\_occ\_change} & \scriptsize yes/no & \footnotesize In the last ten years, did you change your main kind of work?\par\vspace{1pt}{\footnotesize\color{hdr}{\hifont पिछले दस साल में, क्या आपने अपना मुख्य काम बदला?}} & \scriptsize 0 No | 1 Yes & \scriptsize Occupational mobility check; job-history item\\
\footnotesize 33 & \footnotesize\texttt{prev\_occ} & \scriptsize free text & \footnotesize What was your previous main work? Write down exactly what they say, in their words — do not pick from a list.\par\vspace{1pt}{\footnotesize\color{hdr}{\hifont इससे पहले आपका मुख्य काम क्या था?}} & \scriptsize \textit{Asked only if: Ask only if prev\_occ\_change = 1} & \scriptsize Gathmann and Schonberg 2010 (observed moves are systematically shorter than random moves — the validation test for a task-distance measure); coded to NCO-2015 in the office\\
\footnotesize 34 & \footnotesize\texttt{target\_occ} & \scriptsize free text & \footnotesize Looking ahead a few years, what work would you like to be doing? Write down exactly what they say, in their words — do not pick from a list, and do not prompt with examples. If they want to carry on exactly as they are, write that.\par\vspace{1pt}{\footnotesize\color{hdr}{\hifont अगर यह काम बिल्कुल बंद हो जाए, तो आप कौन सा काम करेंगे?}} & \scriptsize  & \scriptsize Project design (stated destination, for the structural-displacement analysis); coded to NCO-2015 in the office\\
\footnotesize 35 & \footnotesize\texttt{prev\_occ\_reason} & \scriptsize select one & \footnotesize What was the main reason you changed?\par\vspace{1pt}{\footnotesize\color{hdr}{\hifont काम बदलने की मुख्य वजह क्या थी?}} & \scriptsize 1 Better income | 2 Lost the previous work | 3 Work ended with the season | 4 Family reasons | 5 Health or injury | 6 Other\par\vspace{1pt}\textit{Asked only if: Ask only if prev\_occ\_change = 1} & \scriptsize Job-history item\\
\bottomrule\end{longtable}\vspace{6pt}
\noindent{\large\bfseries\color{hdr} Module C --- Monthly calendar of work and income}\quad{\small(8 questions)}\\[3pt]
{\small\itshape Read aloud before this module: Now I want to go through the last twelve months, one month at a time. For each month I will ask what work you were mainly doing, and then what you usually earned. It is easiest to start with the Yatra months and fill the rest afterwards.}\\[5pt]
\begin{longtable}{@{}p{0.5cm}p{3.3cm}p{1.7cm}p{7.3cm}p{5.6cm}p{5.4cm}@{}}
\toprule \footnotesize\textbf{\#} & \footnotesize\textbf{Variable} & \footnotesize\textbf{Type} & \footnotesize\textbf{Question (English / Hindi)} & \footnotesize\textbf{Options {\&} skip} & \footnotesize\textbf{Source}\\ \midrule
\endhead
\footnotesize 36 & \footnotesize\texttt{months\_worked\_yatra} & \scriptsize whole number & \footnotesize In the past year, how many months did you work here in the Yatra season? Count the months you spent getting ready for it as well.\par\vspace{1pt}{\footnotesize\color{hdr}{\hifont पिछले साल में आपने यात्रा के सीज़न में यहाँ कितने महीने काम किया? सीज़न की तैयारी के महीने भी गिनें।}} & \scriptsize  & \scriptsize Project design (replaces the twelve-cell status calendar, 2026-10-05)\\
\footnotesize 37 & \footnotesize\texttt{months\_no\_paid\_work} & \scriptsize whole number & \footnotesize And in the past year, how many months did you have no paid work at all?\par\vspace{1pt}{\footnotesize\color{hdr}{\hifont और पिछले साल में कितने महीने आपके पास कोई कमाई वाला काम नहीं था?}} & \scriptsize  & \scriptsize Project design (replaces the twelve-cell status calendar, 2026-10-05)\\
\footnotesize 38 & \footnotesize\texttt{offseason\_activity} & \scriptsize select one & \footnotesize In the remaining months — not the Yatra season, and not the months with no work — what was your main work? Code the closest one.\par\vspace{1pt}{\footnotesize\color{hdr}{\hifont बाकी महीनों में — जो यात्रा के सीज़न के नहीं थे और जिनमें काम भी था — आपका मुख्य काम क्या था? सबसे नज़दीकी चुनें।}} & \scriptsize 2 Farming (crops) | 3 Livestock (animals) | 4 Casual or daily wage labour | 5 Construction work | 6 Own small shop, stall or trade | 7 Salaried job | 96 Some other work | 0 There were no such months\par\vspace{1pt}\textit{Asked only if: May be left blank if the Yatra months and the idle months already account for all twelve} & \scriptsize Project design (replaces the modal non-Yatra calendar cell, 2026-10-05)\\
\footnotesize 39 & \footnotesize\texttt{income\_annual\_total} & \scriptsize amount (Rs) & \footnotesize Thinking of the whole past year, about how much did you earn in total from all your work, after costs? (Rs)\par\vspace{1pt}{\footnotesize\color{hdr}{\hifont पूरे पिछले साल को मिलाकर, सारे काम से ख़र्चा निकालकर आपने कुल कितना कमाया? (रुपये)}} & \scriptsize \textit{Asked only if: May be left blank if the respondent declines to discuss earnings} & \scriptsize Apablaza et al. 2026 Q15 (annual-total fallback), adapted to a seasonal workforce and now the only earnings route\\
\footnotesize 40 & \footnotesize\texttt{pct\_income\_yatra} & \scriptsize whole number & \footnotesize Out of every 100 rupees of that, how many came from your Yatra work? (The rest is counted as coming from your other work.)\par\vspace{1pt}{\footnotesize\color{hdr}{\hifont पूरे साल में आपने जो 100 रुपये कमाए, उनमें से कितने यात्रा के काम से आए? (बाकी दूसरे काम का माना जाएगा।)}} & \scriptsize \textit{Asked only if: Ask only if income\_annual\_total was given; may be left blank} & \scriptsize Project design (the Yatra / non-Yatra split needed to reweight an annual total onto the season)\\
\footnotesize 41 & \footnotesize\texttt{earnings\_range} & \scriptsize select one & \footnotesize Which range is closest to what you usually earn in a month from work? Under Rs 5,000; Rs 5,000 to 10,000; Rs 10,000 to 20,000; Rs 20,000 or more; or don't know?\par\vspace{1pt}{\footnotesize\color{hdr}{\hifont आपकी एक महीने की कमाई लगभग किस दायरे में है? 5,000 रुपये से कम; 5,000 से 10,000; 10,000 से 20,000; 20,000 या उससे ज़्यादा; या पता नहीं?}} & \scriptsize 1 Under Rs 5,000 | 2 Rs 5,000 to 10,000 | 3 Rs 10,000 to 20,000 | 4 Rs 20,000 or more | 97 Don't know\par\vspace{1pt}\textit{Asked only if: Ask only if income\_annual\_total was left blank; may be left blank} & \scriptsize Apablaza et al. 2026, Appendix 2 Q15; bands to check\\
\footnotesize 42 & \footnotesize\texttt{hours\_day\_offseason} & \scriptsize whole number & \footnotesize In the months when the Yatra is closed and you are doing other work, on a normal working day, how many hours do you work in total?\par\vspace{1pt}{\footnotesize\color{hdr}{\hifont जिन महीनों में यात्रा बंद रहती है और आप कोई दूसरा काम करते हैं, उनमें आम दिन में आप कुल कितने घंटे काम करते हैं?}} & \scriptsize \textit{Asked only if: Ask only if there are months that are neither Yatra nor idle (12 - months\_worked\_yatra - months\_no\_paid\_work > 0)} & \scriptsize Apablaza et al. 2026 Q13, asked a second time for the off-season (decomposed as hours/day x days/week, as in the Yatra-season pair)\\
\footnotesize 43 & \footnotesize\texttt{days\_week\_offseason} & \scriptsize whole number & \footnotesize In those months, how many days a week do you usually work?\par\vspace{1pt}{\footnotesize\color{hdr}{\hifont उन महीनों में आप हफ़्ते में कितने दिन काम करते हैं?}} & \scriptsize \textit{Asked only if: Ask only if there are months that are neither Yatra nor idle (12 - months\_worked\_yatra - months\_no\_paid\_work > 0)} & \scriptsize Apablaza et al. 2026 Q13, asked a second time for the off-season\\
\bottomrule\end{longtable}\vspace{6pt}
\noindent{\large\bfseries\color{hdr} Module D --- Migration, home place and remittances}\quad{\small(17 questions)}\\[3pt]
{\small\itshape Read aloud before this module: Now a few questions about where your home is, and what you do when the Yatra closes for the season.}\\[5pt]
\begin{longtable}{@{}p{0.5cm}p{3.3cm}p{1.7cm}p{7.3cm}p{5.6cm}p{5.4cm}@{}}
\toprule \footnotesize\textbf{\#} & \footnotesize\textbf{Variable} & \footnotesize\textbf{Type} & \footnotesize\textbf{Question (English / Hindi)} & \footnotesize\textbf{Options {\&} skip} & \footnotesize\textbf{Source}\\ \midrule
\endhead
\footnotesize 44 & \footnotesize\texttt{resp\_returns\_at\_closure} & \scriptsize yes/no & \footnotesize When the Yatra closes for the season, do YOU go to your home place?\par\vspace{1pt}{\footnotesize\color{hdr}{\hifont जब यात्रा का मौसम ख़त्म हो जाता है, तो क्या आप अपने घर चले जाते हैं?}} & \scriptsize 0 No | 1 Yes & \scriptsize Project design; categories from the pilot's own residency\_pattern distribution (n=46)\\
\footnotesize 45 & \footnotesize\texttt{hh\_at\_home\_place} & \scriptsize yes/no & \footnotesize For most of the year, do the rest of your household live in your home village or town?\par\vspace{1pt}{\footnotesize\color{hdr}{\hifont साल के ज़्यादातर समय, क्या आपके घर के बाक़ी लोग आपके गाँव या शहर में ही रहते हैं?}} & \scriptsize 0 No | 1 Yes & \scriptsize Project design (split-household measurement, underpinning Module E)\\
\footnotesize 46 & \footnotesize\texttt{remit\_differs\_by\_season} & \scriptsize yes/no & \footnotesize Is the money you send home different in the Yatra season from when the Yatra is closed?\par\vspace{1pt}{\footnotesize\color{hdr}{\hifont यात्रा के सीज़न में आप घर जो पैसा भेजते हैं, वह यात्रा बंद रहने के महीनों से अलग होता है?}} & \scriptsize 0 No | 1 Yes & \scriptsize Project design (mirrors spend\_differs\_by\_season, Module E)\\
\footnotesize 47 & \footnotesize\texttt{remit\_out\_yatra\_pm} & \scriptsize amount (Rs) & \footnotesize In a normal month during the Yatra season, how much money do you usually send to family or others living elsewhere? (Rs, 0 if none)\par\vspace{1pt}{\footnotesize\color{hdr}{\hifont यात्रा के मौसम के आम महीने में, आप घर या कहीं और रहने वालों को कितना पैसा भेजते हैं? (रुपये, न भेजते हों तो 0)}} & \scriptsize  & \scriptsize Project design (amount only; frequency dropped for length). A previous citation to NSS 64th Round practice was withdrawn 2026-09-28: no NSS schedule is held by this project and it could not be checked\\
\footnotesize 48 & \footnotesize\texttt{remit\_out\_offseason\_pm} & \scriptsize amount (Rs) & \footnotesize In a normal month when the Yatra is closed, how much do you usually send? (Rs, 0 if none)\par\vspace{1pt}{\footnotesize\color{hdr}{\hifont जब यात्रा बंद रहती है, उस आम महीने में कितना भेजते हैं? (रुपये, न भेजते हों तो 0)}} & \scriptsize \textit{Asked only if: Ask only if remit\_differs\_by\_season = 1} & \scriptsize Project design (amount only; frequency dropped for length). A previous citation to NSS 64th Round practice was withdrawn 2026-09-28: no NSS schedule is held by this project and it could not be checked\\
\footnotesize 49 & \footnotesize\texttt{remit\_mode} & \scriptsize select one & \footnotesize How do you usually send it?\par\vspace{1pt}{\footnotesize\color{hdr}{\hifont आप पैसा आम तौर पर कैसे भेजते हैं?}} & \scriptsize 1 From my own phone (UPI or net banking) | 2 In person at a bank branch or bank mitra / CSP | 3 Money order or post office | 4 Sent with someone going home | 5 Carried it myself | 6 Through an agent or middleman | 7 Other\par\vspace{1pt}\textit{Asked only if: Ask only if remit\_out\_yatra\_pm > 0 or remit\_out\_offseason\_pm > 0} & \scriptsize Project design (remittance channel as a financial-inclusion and cost measure)\\
\footnotesize 50 & \footnotesize\texttt{remit\_in\_yatra\_pm} & \scriptsize amount (Rs) & \footnotesize In a normal month during the Yatra season, how much money do you usually receive from family or others? (Rs, 0 if none)\par\vspace{1pt}{\footnotesize\color{hdr}{\hifont यात्रा के मौसम के आम महीने में, घरवालों या किसी और से आपको कितना पैसा मिलता है? (रुपये, न मिलता हो तो 0)}} & \scriptsize  & \scriptsize Project design (amount only; frequency dropped for length). A previous citation to NSS 64th Round practice was withdrawn 2026-09-28: no NSS schedule is held by this project and it could not be checked\\
\footnotesize 51 & \footnotesize\texttt{remit\_in\_offseason\_pm} & \scriptsize amount (Rs) & \footnotesize In a normal month when the Yatra is closed, how much do you usually receive? (Rs, 0 if none)\par\vspace{1pt}{\footnotesize\color{hdr}{\hifont जब यात्रा बंद रहती है, उस आम महीने में कितना मिलता है? (रुपये, न मिलता हो तो 0)}} & \scriptsize \textit{Asked only if: Ask only if remit\_differs\_by\_season = 1} & \scriptsize Project design (amount only; frequency dropped for length). A previous citation to NSS 64th Round practice was withdrawn 2026-09-28: no NSS schedule is held by this project and it could not be checked\\
\footnotesize 52 & \footnotesize\texttt{n\_hh\_same\_business} & \scriptsize whole number & \footnotesize How many people from your household work in the same business or establishment as you, counting yourself?\par\vspace{1pt}{\footnotesize\color{hdr}{\hifont आपके घर के कितने लोग आपके साथ उसी काम या उसी दुकान/धंधे में काम करते हैं, आपको भी गिनकर?}} & \scriptsize  & \scriptsize Project design (enterprise concentration of household labour, 2026-10-05)\\
\footnotesize 53 & \footnotesize\texttt{years\_coming\_here} & \scriptsize whole number & \footnotesize How many years have you been coming here for the Yatra season?\par\vspace{1pt}{\footnotesize\color{hdr}{\hifont आप कितने सालों से यात्रा के मौसम में यहाँ आ रहे हैं?}} & \scriptsize \textit{Asked only if: Ask only if resp\_returns\_at\_closure = 1} & \scriptsize Project design (duration of the seasonal relationship)\\
\footnotesize 54 & \footnotesize\texttt{came\_here\_reason} & \scriptsize select one & \footnotesize What was the main reason you first came here to work?\par\vspace{1pt}{\footnotesize\color{hdr}{\hifont आप सबसे पहले यहाँ काम करने किस मुख्य वजह से आए थे?}} & \scriptsize 1 No work at home | 2 Pay is better here | 3 Family or people from my village were already here | 4 Land at home is too little to live on | 5 Debt to repay | 6 A contractor or agent brought me | 7 Married into / moved with family | 8 Some other reason\par\vspace{1pt}\textit{Asked only if: Ask only if origin is not Local (same district)} & \scriptsize Project design (push/pull as a vulnerability covariate)\\
\footnotesize 55 & \footnotesize\texttt{came\_here\_reason\_other} & \scriptsize free text & \footnotesize What was that reason? Write down what they say.\par\vspace{1pt}{\footnotesize\color{hdr}{\hifont वह वजह क्या थी? जो वे कहें वही लिख लें।}} & \scriptsize \textit{Asked only if: Ask only if came\_here\_reason = 8 (some other reason). May be left blank.} & \scriptsize Project design (catch-all codes lose the case they catch)\\
\footnotesize 56 & \footnotesize\texttt{worked\_away\_in\_closure} & \scriptsize yes/no & \footnotesize Last year, in the months when the Yatra was closed, did you work away from your home place?\par\vspace{1pt}{\footnotesize\color{hdr}{\hifont पिछले साल, जिन महीनों में यात्रा बंद थी, क्या आपने अपने घर वाली जगह से बाहर जाकर काम किया?}} & \scriptsize 0 No | 1 Yes & \scriptsize Project design (off-season labour migration)\\
\footnotesize 57 & \footnotesize\texttt{closure\_work\_detail} & \scriptsize free text & \footnotesize Where did you go, and what work did you do there? Write what they say.\par\vspace{1pt}{\footnotesize\color{hdr}{\hifont कहाँ गए थे, और वहाँ क्या काम किया? जैसा वे बताएँ वैसा लिख लीजिए।}} & \scriptsize \textit{Asked only if: Ask only if worked\_away\_in\_closure = 1; may be left blank} & \scriptsize Project design; coded to NCO-2015 in the office\\
\footnotesize 58 & \footnotesize\texttt{would\_move\_for\_work} & \scriptsize select one & \footnotesize If this work here ended, would you go away from your home place to look for work in the coming year?\par\vspace{1pt}{\footnotesize\color{hdr}{\hifont अगर यहाँ का यह काम बंद हो जाए, तो क्या आप आने वाले साल में काम ढूँढने अपने घर से बाहर जाएँगे?}} & \scriptsize 0 No | 1 Yes | 97 Don't know & \scriptsize Project design; horizoned intention item after VASyR 2025 (move\_accom\_yesno, asked over a stated 6-month horizon)\\
\footnotesize 59 & \footnotesize\texttt{migration\_referral} & \scriptsize select one & \footnotesize Who mainly helped you get this work, or arranged it for you?\par\vspace{1pt}{\footnotesize\color{hdr}{\hifont यहाँ आने या यह काम पकड़ने में मुख्य रूप से किसने मदद की या रास्ता दिखाया: यहाँ पहले से रह रहा कोई घरवाला; कोई दोस्त या जान-पहचान वाला; कोई ठेकेदार या एजेंट; कोई नेता या समाज का बड़ा आदमी; कोई नहीं, आपने खुद ढूँढा; मालिक ने सीधे बुलाया; या कोई और?}} & \scriptsize 1 A family member already working here | 2 A friend or someone from the village | 3 A thekedar or contractor I now work for | 4 An agent or middleman who placed me (usually for a fee) | 5 No one; I found it myself | 6 The employer called me directly | 7 Someone else & \scriptsize Project design (referral channel as a proxy for social capital and for placement debt)\\
\footnotesize 60 & \footnotesize\texttt{migration\_referral\_other} & \scriptsize free text & \footnotesize Who was it? Write what they say.\par\vspace{1pt}{\footnotesize\color{hdr}{\hifont वे कौन थे? जैसा वे बताएँ वैसा लिख लीजिए।}} & \scriptsize \textit{Asked only if: Ask only if migration\_referral = 7 (Other); may be left blank} & \scriptsize Standard household-survey item\\
\bottomrule\end{longtable}\vspace{6pt}
\noindent{\large\bfseries\color{hdr} Module E --- Household consumption}\quad{\small(27 questions)}\\[3pt]
{\small\itshape Read aloud before this module: Now about what your household usually spends in a month. For each thing I will ask twice — once for a normal month during the Yatra season, and once for a normal month when the Yatra is closed — because spending is often quite different in the two. Rough figures are fine.}\\[5pt]
\begin{longtable}{@{}p{0.5cm}p{3.3cm}p{1.7cm}p{7.3cm}p{5.6cm}p{5.4cm}@{}}
\toprule \footnotesize\textbf{\#} & \footnotesize\textbf{Variable} & \footnotesize\textbf{Type} & \footnotesize\textbf{Question (English / Hindi)} & \footnotesize\textbf{Options {\&} skip} & \footnotesize\textbf{Source}\\ \midrule
\endhead
\footnotesize 61 & \footnotesize\texttt{spend\_differs\_by\_season} & \scriptsize yes/no & \footnotesize Leaving aside money you send home — is what your household spends in a normal month during the Yatra season different from what it spends when the Yatra is closed?\par\vspace{1pt}{\footnotesize\color{hdr}{\hifont घर भेजे जाने वाले पैसे को छोड़कर — यात्रा के मौसम में आपके घर का महीने का ख़र्च, और जब यात्रा बंद रहती है तब का ख़र्च, क्या अलग-अलग होता है?}} & \scriptsize 0 No | 1 Yes & \scriptsize Project design (respondent-gated seasonal recall); the underlying two-season design follows Beegle et al. 2012 and Deaton and Grosh 2000 on usual-period recall\\
\footnotesize 62 & \footnotesize\texttt{cons\_staples\_yatra\_pm} & \scriptsize amount (Rs) & \footnotesize In a normal month during the Yatra season, how much does your household spend on cereals (rice, wheat, other grains), pulses, sugar and salt, bought or from your own stock? (Rs)\par\vspace{1pt}{\footnotesize\color{hdr}{\hifont यात्रा के मौसम के आम महीने में, आप (और जो यहाँ आपके साथ रहते हैं) अनाज (चावल, गेहूँ, दूसरा अनाज), दाल, चीनी और नमक पर, ख़रीदा हुआ या अपने घर का कितना ख़र्च करते हैं? (रुपये)}} & \scriptsize  & \scriptsize HCES 2022-23, MoSPI Appendix A, Sections 5.1-5.3 (30-day); IHDS-II Income and Social Capital Questionnaire, Q14 14.1-14.6 (30-day, same items)\\
\footnotesize 63 & \footnotesize\texttt{cons\_staples\_offseason\_pm} & \scriptsize amount (Rs) & \footnotesize In a normal month when the Yatra is closed, how much does your household spend on cereals (rice, wheat, other grains), pulses, sugar and salt, bought or from your own stock? (Rs)\par\vspace{1pt}{\footnotesize\color{hdr}{\hifont जब यात्रा बंद रहती है, उस आम महीने में आपका घर अनाज (चावल, गेहूँ, दूसरा अनाज), दाल, चीनी और नमक पर, ख़रीदा हुआ या अपने घर का कितना ख़र्च करता है? (रुपये)}} & \scriptsize \textit{Asked only if: Ask only if spend\_differs\_by\_season = 1} & \scriptsize HCES 2022-23, MoSPI Appendix A, Sections 5.1-5.3 (30-day); IHDS-II Income and Social Capital Questionnaire, Q14 14.1-14.6 (30-day, same items)\\
\footnotesize 64 & \footnotesize\texttt{cons\_perishables\_yatra\_pm} & \scriptsize amount (Rs) & \footnotesize In a normal month during the Yatra season, how much does your household spend on milk and milk products, vegetables, fruit, egg/fish/meat, cooking oil, spices, and tea or coffee, bought or from your own stock? (Rs)\par\vspace{1pt}{\footnotesize\color{hdr}{\hifont यात्रा के मौसम के आम महीने में, आप (और जो यहाँ आपके साथ रहते हैं) दूध और दूध की चीज़ें, सब्ज़ी, फल, अंडा/मछली/मांस, तेल, मसाले, और चाय पर कितना ख़र्च करते हैं? (रुपये)}} & \scriptsize  & \scriptsize HCES 2022-23, MoSPI Appendix A, Sections 6.1-6.8 (7-day actual recall in HCES; asked here as a usual monthly figure instead, for the reason above)\\
\footnotesize 65 & \footnotesize\texttt{cons\_perishables\_offseason\_pm} & \scriptsize amount (Rs) & \footnotesize In a normal month when the Yatra is closed, how much does your household spend on milk and milk products, vegetables, fruit, egg/fish/meat, cooking oil, spices, and tea or coffee, bought or from your own stock? (Rs)\par\vspace{1pt}{\footnotesize\color{hdr}{\hifont जब यात्रा बंद रहती है, उस आम महीने में आपका घर दूध और दूध की चीज़ें, सब्ज़ी, फल, अंडा/मछली/मांस, तेल, मसाले, और चाय पर कितना ख़र्च करता है? (रुपये)}} & \scriptsize \textit{Asked only if: Ask only if spend\_differs\_by\_season = 1} & \scriptsize HCES 2022-23, MoSPI Appendix A, Sections 6.1-6.8 (7-day actual recall in HCES; asked here as a usual monthly figure instead, for the reason above)\\
\footnotesize 66 & \footnotesize\texttt{cons\_food\_own\_yatra\_pm} & \scriptsize amount (Rs) & \footnotesize If your household had not grown, raised or been given any of your food, about how much would it have cost to buy in a normal month during the Yatra season? (Rs, 0 if you buy everything)\par\vspace{1pt}{\footnotesize\color{hdr}{\hifont अगर आप (और जो यहाँ आपके साथ रहते हैं) अपना खाना ख़ुद न उगाते, न पालते और न किसी से मिलता, तो यात्रा के मौसम के आम महीने में उसे ख़रीदने में लगभग कितना लगता? (रुपये, सब ख़रीदते हों तो 0)}} & \scriptsize  & \scriptsize Calvo and Dercon 2007; Eze and Iheonu 2025 (non-purchased food, aggregate approach)\\
\footnotesize 67 & \footnotesize\texttt{cons\_food\_own\_offseason\_pm} & \scriptsize amount (Rs) & \footnotesize If your household had not grown, raised or been given any of your food, about how much would it have cost to buy in a normal month when the Yatra is closed? (Rs, 0 if you buy everything)\par\vspace{1pt}{\footnotesize\color{hdr}{\hifont अगर आपका घर अपना खाना ख़ुद न उगाता, न पालता और न किसी से मिलता, तो जब यात्रा बंद रहती है उस आम महीने में उसे ख़रीदने में लगभग कितना लगता? (रुपये, सब ख़रीदते हों तो 0)}} & \scriptsize \textit{Asked only if: Ask only if spend\_differs\_by\_season = 1} & \scriptsize Calvo and Dercon 2007; Eze and Iheonu 2025 (non-purchased food, aggregate approach)\\
\footnotesize 68 & \footnotesize\texttt{cons\_food\_out\_yatra\_pm} & \scriptsize amount (Rs) & \footnotesize In a normal month during the Yatra season, how much does your household spend on meals, tea and snacks eaten outside the home, including any an employer gave free, at their market value? (Rs)\par\vspace{1pt}{\footnotesize\color{hdr}{\hifont यात्रा के मौसम के आम महीने में, आप (और जो यहाँ आपके साथ रहते हैं) बाहर खाए गए खाने, चाय और नाश्ते पर, मालिक ने मुफ़्त दिया हो तो उसकी क़ीमत मिलाकर कितना ख़र्च करते हैं? (रुपये)}} & \scriptsize  & \scriptsize HCES 2022-23, MoSPI Appendix A, Section 7.1 'served processed food' (7-day in HCES; usual monthly here)\\
\footnotesize 69 & \footnotesize\texttt{cons\_food\_out\_offseason\_pm} & \scriptsize amount (Rs) & \footnotesize In a normal month when the Yatra is closed, how much does your household spend on meals, tea and snacks eaten outside the home, including any an employer gave free, at their market value? (Rs)\par\vspace{1pt}{\footnotesize\color{hdr}{\hifont जब यात्रा बंद रहती है, उस आम महीने में आपका घर बाहर खाए गए खाने, चाय और नाश्ते पर, मालिक ने मुफ़्त दिया हो तो उसकी क़ीमत मिलाकर कितना ख़र्च करता है? (रुपये)}} & \scriptsize \textit{Asked only if: Ask only if spend\_differs\_by\_season = 1} & \scriptsize HCES 2022-23, MoSPI Appendix A, Section 7.1 'served processed food' (7-day in HCES; usual monthly here)\\
\footnotesize 70 & \footnotesize\texttt{cons\_fuel\_yatra\_pm} & \scriptsize amount (Rs) & \footnotesize In a normal month during the Yatra season, how much does your household spend on fuel and light: cooking fuel, firewood, electricity, kerosene, candles? (Rs)\par\vspace{1pt}{\footnotesize\color{hdr}{\hifont यात्रा के मौसम के आम महीने में, आप (और जो यहाँ आपके साथ रहते हैं) ईंधन और रोशनी पर: खाना बनाने की गैस, लकड़ी, बिजली, मिट्टी का तेल, मोमबत्ती कितना ख़र्च करते हैं? (रुपये)}} & \scriptsize  & \scriptsize HCES 2022-23, MoSPI Appendix A, Section 8.1 (30-day). IHDS-II Income and Social Capital Questionnaire, Q14 14.21-14.22. Nigeria GHS-Panel Wave 3, Household Questionnaire, Sections 10B/11 (30-day tier)\\
\footnotesize 71 & \footnotesize\texttt{cons\_fuel\_offseason\_pm} & \scriptsize amount (Rs) & \footnotesize In a normal month when the Yatra is closed, how much does your household spend on fuel and light: cooking fuel, firewood, electricity, kerosene, candles? (Rs)\par\vspace{1pt}{\footnotesize\color{hdr}{\hifont जब यात्रा बंद रहती है, उस आम महीने में आपका घर ईंधन और रोशनी पर: खाना बनाने की गैस, लकड़ी, बिजली, मिट्टी का तेल, मोमबत्ती कितना ख़र्च करता है? (रुपये)}} & \scriptsize \textit{Asked only if: Ask only if spend\_differs\_by\_season = 1} & \scriptsize HCES 2022-23, MoSPI Appendix A, Section 8.1 (30-day). IHDS-II Income and Social Capital Questionnaire, Q14 14.21-14.22. Nigeria GHS-Panel Wave 3, Household Questionnaire, Sections 10B/11 (30-day tier)\\
\footnotesize 72 & \footnotesize\texttt{cons\_routine\_misc\_yatra\_pm} & \scriptsize amount (Rs) & \footnotesize In a normal month during the Yatra season, how much does your household spend on soap, toiletries, cleaning goods and other small household items? (Rs)\par\vspace{1pt}{\footnotesize\color{hdr}{\hifont यात्रा के मौसम के आम महीने में, आप (और जो यहाँ आपके साथ रहते हैं) साबुन, नहाने-धोने का सामान, सफ़ाई का सामान और घर की दूसरी छोटी चीज़ों पर कितना ख़र्च करते हैं? (रुपये)}} & \scriptsize  & \scriptsize HCES 2022-23, MoSPI Appendix A, Sections 9.1-9.2 (30-day). IHDS-II Income and Social Capital Questionnaire, Q14 14.25-14.27. Nigeria GHS-Panel Wave 3, Household Questionnaire, Sections 10B/11 (30-day tier)\\
\footnotesize 73 & \footnotesize\texttt{cons\_routine\_misc\_offseason\_pm} & \scriptsize amount (Rs) & \footnotesize In a normal month when the Yatra is closed, how much does your household spend on soap, toiletries, cleaning goods and other small household items? (Rs)\par\vspace{1pt}{\footnotesize\color{hdr}{\hifont जब यात्रा बंद रहती है, उस आम महीने में आपका घर साबुन, नहाने-धोने का सामान, सफ़ाई का सामान और घर की दूसरी छोटी चीज़ों पर कितना ख़र्च करता है? (रुपये)}} & \scriptsize \textit{Asked only if: Ask only if spend\_differs\_by\_season = 1} & \scriptsize HCES 2022-23, MoSPI Appendix A, Sections 9.1-9.2 (30-day). IHDS-II Income and Social Capital Questionnaire, Q14 14.25-14.27. Nigeria GHS-Panel Wave 3, Household Questionnaire, Sections 10B/11 (30-day tier)\\
\footnotesize 74 & \footnotesize\texttt{cons\_transport\_comm\_yatra\_pm} & \scriptsize amount (Rs) & \footnotesize In a normal month during the Yatra season, how much does your household spend on local transport (bus, shared jeep, auto) and phone or internet charges? (Rs)\par\vspace{1pt}{\footnotesize\color{hdr}{\hifont यात्रा के मौसम के आम महीने में, आप (और जो यहाँ आपके साथ रहते हैं) आने-जाने (बस, जीप, ऑटो) और फ़ोन या इंटरनेट पर कितना ख़र्च करते हैं? (रुपये)}} & \scriptsize  & \scriptsize HCES 2022-23, MoSPI Appendix A, Sections 11.1-11.2 (30-day). IHDS-II Income and Social Capital Questionnaire, Q14 14.24/14.28. Nigeria GHS-Panel Wave 3, Household Questionnaire, Sections 10B/11 (30-day tier)\\
\footnotesize 75 & \footnotesize\texttt{cons\_transport\_comm\_offseason\_pm} & \scriptsize amount (Rs) & \footnotesize In a normal month when the Yatra is closed, how much does your household spend on local transport (bus, shared jeep, auto) and phone or internet charges? (Rs)\par\vspace{1pt}{\footnotesize\color{hdr}{\hifont जब यात्रा बंद रहती है, उस आम महीने में आपका घर आने-जाने (बस, जीप, ऑटो) और फ़ोन या इंटरनेट पर कितना ख़र्च करता है? (रुपये)}} & \scriptsize \textit{Asked only if: Ask only if spend\_differs\_by\_season = 1} & \scriptsize HCES 2022-23, MoSPI Appendix A, Sections 11.1-11.2 (30-day). IHDS-II Income and Social Capital Questionnaire, Q14 14.24/14.28. Nigeria GHS-Panel Wave 3, Household Questionnaire, Sections 10B/11 (30-day tier)\\
\footnotesize 76 & \footnotesize\texttt{cons\_rent\_yatra\_pm} & \scriptsize amount (Rs) & \footnotesize In a normal month during the Yatra season, how much does your household spend on rent for the home (0 if owned, and 0 for any period spent living for free)? (Rs)\par\vspace{1pt}{\footnotesize\color{hdr}{\hifont यात्रा के मौसम के आम महीने में, आप (और जो यहाँ आपके साथ रहते हैं) घर के किराये पर (अपना घर हो तो 0, और मुफ़्त रहे हों तो भी 0) कितना ख़र्च करते हैं? (रुपये)}} & \scriptsize  & \scriptsize HCES 2022-23, MoSPI Appendix A, Section 11.4 (30-day). IHDS-II Income and Social Capital Questionnaire, Q14 14.30. Nigeria GHS-Panel Wave 3, Household Questionnaire, Sections 10B/11 (30-day tier)\\
\footnotesize 77 & \footnotesize\texttt{cons\_rent\_offseason\_pm} & \scriptsize amount (Rs) & \footnotesize In a normal month when the Yatra is closed, how much does your household spend on rent for the home (0 if owned, and 0 for any period spent living for free)? (Rs)\par\vspace{1pt}{\footnotesize\color{hdr}{\hifont जब यात्रा बंद रहती है, उस आम महीने में आपका घर घर के किराये पर (अपना घर हो तो 0, और मुफ़्त रहे हों तो भी 0) कितना ख़र्च करता है? (रुपये)}} & \scriptsize \textit{Asked only if: Ask only if spend\_differs\_by\_season = 1} & \scriptsize HCES 2022-23, MoSPI Appendix A, Section 11.4 (30-day). IHDS-II Income and Social Capital Questionnaire, Q14 14.30. Nigeria GHS-Panel Wave 3, Household Questionnaire, Sections 10B/11 (30-day tier)\\
\footnotesize 78 & \footnotesize\texttt{cons\_med\_nonhosp\_yatra\_pm} & \scriptsize amount (Rs) & \footnotesize In a normal month during the Yatra season, how much does your household spend on medicine, doctor's or clinic fees and tests, not counting a hospital stay? (Rs)\par\vspace{1pt}{\footnotesize\color{hdr}{\hifont यात्रा के मौसम के आम महीने में, आप (और जो यहाँ आपके साथ रहते हैं) दवाई, डॉक्टर या क्लीनिक की फ़ीस और जाँच पर, अस्पताल में भर्ती को छोड़कर कितना ख़र्च करते हैं? (रुपये)}} & \scriptsize  & \scriptsize HCES 2022-23, MoSPI Appendix A, Section 10.3 (30-day). IHDS-II Income and Social Capital Questionnaire, Q14 14.33 (30-day, out-patient)\\
\footnotesize 79 & \footnotesize\texttt{cons\_med\_nonhosp\_offseason\_pm} & \scriptsize amount (Rs) & \footnotesize In a normal month when the Yatra is closed, how much does your household spend on medicine, doctor's or clinic fees and tests, not counting a hospital stay? (Rs)\par\vspace{1pt}{\footnotesize\color{hdr}{\hifont जब यात्रा बंद रहती है, उस आम महीने में आपका घर दवाई, डॉक्टर या क्लीनिक की फ़ीस और जाँच पर, अस्पताल में भर्ती को छोड़कर कितना ख़र्च करता है? (रुपये)}} & \scriptsize \textit{Asked only if: Ask only if spend\_differs\_by\_season = 1} & \scriptsize HCES 2022-23, MoSPI Appendix A, Section 10.3 (30-day). IHDS-II Income and Social Capital Questionnaire, Q14 14.33 (30-day, out-patient)\\
\footnotesize 80 & \footnotesize\texttt{cons\_packaged\_food\_yatra\_pm} & \scriptsize amount (Rs) & \footnotesize In a normal month during the Yatra season, how much does your household spend on biscuits, namkeen, chips, packaged snacks, cold drinks or bottled water? (Rs)\par\vspace{1pt}{\footnotesize\color{hdr}{\hifont यात्रा के मौसम के आम महीने में, आप (और जो यहाँ आपके साथ रहते हैं) बिस्किट, नमकीन, चिप्स, पैकेट वाले नाश्ते, कोल्ड ड्रिंक या बोतलबंद पानी पर कितना ख़र्च करते हैं? (रुपये)}} & \scriptsize  & \scriptsize HCES 2022-23, MoSPI Appendix A, Section 7.2 (7-day in HCES; usual monthly here, as for every other seasonal item)\\
\footnotesize 81 & \footnotesize\texttt{cons\_packaged\_food\_offseason\_pm} & \scriptsize amount (Rs) & \footnotesize In a normal month when the Yatra is closed, how much does your household spend on biscuits, namkeen, chips, packaged snacks, cold drinks or bottled water? (Rs)\par\vspace{1pt}{\footnotesize\color{hdr}{\hifont जब यात्रा बंद रहती है, उस आम महीने में आपका घर बिस्किट, नमकीन, चिप्स, पैकेट वाले नाश्ते, कोल्ड ड्रिंक या बोतलबंद पानी पर कितना ख़र्च करता है? (रुपये)}} & \scriptsize \textit{Asked only if: Ask only if spend\_differs\_by\_season = 1} & \scriptsize HCES 2022-23, MoSPI Appendix A, Section 7.2 (7-day in HCES; usual monthly here, as for every other seasonal item)\\
\footnotesize 82 & \footnotesize\texttt{cons\_pan\_tobacco\_yatra\_pm} & \scriptsize amount (Rs) & \footnotesize In a normal month during the Yatra season, how much does your household spend on pan, gutka, bidi, cigarettes, tobacco or alcohol? (Rs)\par\vspace{1pt}{\footnotesize\color{hdr}{\hifont यात्रा के मौसम के आम महीने में, आप (और जो यहाँ आपके साथ रहते हैं) पान, गुटका, बीड़ी, सिगरेट, तंबाकू या शराब पर कितना ख़र्च करते हैं? (रुपये)}} & \scriptsize  & \scriptsize HCES 2022-23, MoSPI Appendix A, Sections 12.1-12.3 (7-day in HCES; usual monthly here, and combined into one figure)\\
\footnotesize 83 & \footnotesize\texttt{cons\_pan\_tobacco\_offseason\_pm} & \scriptsize amount (Rs) & \footnotesize In a normal month when the Yatra is closed, how much does your household spend on pan, gutka, bidi, cigarettes, tobacco or alcohol? (Rs)\par\vspace{1pt}{\footnotesize\color{hdr}{\hifont जब यात्रा बंद रहती है, उस आम महीने में आपका घर पान, गुटका, बीड़ी, सिगरेट, तंबाकू या शराब पर कितना ख़र्च करता है? (रुपये)}} & \scriptsize \textit{Asked only if: Ask only if spend\_differs\_by\_season = 1} & \scriptsize HCES 2022-23, MoSPI Appendix A, Sections 12.1-12.3 (7-day in HCES; usual monthly here, and combined into one figure)\\
\footnotesize 84 & \footnotesize\texttt{cons\_clothing\_12m} & \scriptsize amount (Rs) & \footnotesize In the last 12 months, how much did your household spend on clothes and footwear? (Rs)\par\vspace{1pt}{\footnotesize\color{hdr}{\hifont पिछले 12 महीनों में आपके घर ने कपड़ों और जूते-चप्पल पर कितना ख़र्च किया? (रुपये)}} & \scriptsize  & \scriptsize HCES 2022-23, MoSPI Appendix A, Sections 13.1-13.2 (365-day). IHDS-II Income and Social Capital Questionnaire, Q14 14.38-14.39 (365-day). Nigeria GHS-Panel Wave 3, Household Questionnaire, Sections 10B/11 instead uses a 6-month recall for clothing (items 401-441); we follow the two India sources, since one of them (HCES) underlies our poverty line\\
\footnotesize 85 & \footnotesize\texttt{cons\_education\_12m} & \scriptsize amount (Rs) & \footnotesize In the last 12 months, how much on education: fees, books, uniforms, tuition? (Rs)\par\vspace{1pt}{\footnotesize\color{hdr}{\hifont पिछले 12 महीनों में पढ़ाई पर कितना ख़र्च हुआ: फ़ीस, किताबें, ड्रेस, ट्यूशन? (रुपये)}} & \scriptsize  & \scriptsize HCES 2022-23, MoSPI Appendix A, Section 10.1 (365-day). IHDS-II Income and Social Capital Questionnaire, Q14 14.35-14.37 (365-day, same items). Nigeria GHS-Panel Wave 3, Household Questionnaire, Sections 10B/11 asks school fees inside its education/roster module, per child, not as a household total, so it is not directly comparable\\
\footnotesize 86 & \footnotesize\texttt{cons\_medical\_hosp\_12m} & \scriptsize amount (Rs) & \footnotesize In the last 12 months, how much did your household pay for any hospital admission or stay? (Rs, 0 if none)\par\vspace{1pt}{\footnotesize\color{hdr}{\hifont पिछले 12 महीनों में अस्पताल में भर्ती होने पर आपके घर का कितना ख़र्च हुआ? (रुपये, न हुआ हो तो 0)}} & \scriptsize  & \scriptsize HCES 2022-23, MoSPI Appendix A, Section 10.2 (365-day, hospitalisation). IHDS-II Income and Social Capital Questionnaire, Q14 14.34 (365-day, in-patient). Nigeria GHS-Panel Wave 3, Household Questionnaire, Sections 10B/11 instead uses one combined 6-month medical item (less detail); we follow the matching HCES/IHDS split\\
\footnotesize 87 & \footnotesize\texttt{cons\_durables\_12m} & \scriptsize amount (Rs) & \footnotesize In the last 12 months, how much on durable goods: furniture, utensils, cooking appliances, phone, jewellery or ornaments, bicycle or vehicle parts? (Rs)\par\vspace{1pt}{\footnotesize\color{hdr}{\hifont पिछले 12 महीनों में सामान पर कितना ख़र्च हुआ: फ़र्नीचर, बर्तन, चूल्हा-मशीन, फ़ोन, गहने, साइकिल या गाड़ी के पुर्ज़े? (रुपये)}} & \scriptsize  & \scriptsize HCES 2022-23, MoSPI Appendix A, Section 14 (365-day, 10 sub-groups). IHDS-II Income and Social Capital Questionnaire, Q14 14.40-14.48 (365-day). Nigeria GHS-Panel Wave 3, Household Questionnaire, Sections 10B/11 (12-month tier: durables, appliances, building materials). All three agree on the 12-month recall; insurance premiums, which IHDS (14.50) and the Nigeria survey (items 510-513) count here, are excluded, following HCES and standard national-accounts practice (insurance is a financial item, not consumption) — premiums are covered instead by the insurance items in Module G\\
\bottomrule\end{longtable}\vspace{6pt}
\noindent{\large\bfseries\color{hdr} Module F --- Housing, amenities and assets}\quad{\small(17 questions)}\\[3pt]
{\small\itshape Read aloud before this module: Now some questions about your house and the things your household has.}\\[5pt]
\begin{longtable}{@{}p{0.5cm}p{3.3cm}p{1.7cm}p{7.3cm}p{5.6cm}p{5.4cm}@{}}
\toprule \footnotesize\textbf{\#} & \footnotesize\textbf{Variable} & \footnotesize\textbf{Type} & \footnotesize\textbf{Question (English / Hindi)} & \footnotesize\textbf{Options {\&} skip} & \footnotesize\textbf{Source}\\ \midrule
\endhead
\footnotesize 88 & \footnotesize\texttt{floor\_material} & \scriptsize select one & \footnotesize What is the main material of the floor at your usual home?\par\vspace{1pt}{\footnotesize\color{hdr}{\hifont आपके असली घर में फ़र्श मुख्य रूप से किस चीज़ का है?}} & \scriptsize 1 Mud/kaccha | 2 Cement/mud-cement | 3 Tile/mosaic/marble & \scriptsize NITI Aayog National MPI housing indicator\\
\footnotesize 89 & \footnotesize\texttt{roof\_material} & \scriptsize select one & \footnotesize What is the main material of the roof at your usual home?\par\vspace{1pt}{\footnotesize\color{hdr}{\hifont आपके असली घर की छत मुख्य रूप से किस चीज़ की है?}} & \scriptsize 1 Thatch/wood/mud | 2 Tin/GI sheet | 3 Concrete/RCC & \scriptsize NITI Aayog National MPI housing indicator\\
\footnotesize 90 & \footnotesize\texttt{wall\_material} & \scriptsize select one & \footnotesize What are the walls of your usual home mainly made of?\par\vspace{1pt}{\footnotesize\color{hdr}{\hifont आपके असली घर की दीवारें मुख्य रूप से किस चीज़ की बनी हैं?}} & \scriptsize 1 Mud, thatch, bamboo or other natural material | 2 Unburnt brick, wood or tin | 3 Burnt brick, cement, concrete or stone & \scriptsize NITI Aayog National MPI housing indicator (wall limb)\\
\footnotesize 91 & \footnotesize\texttt{accom\_type\_here} & \scriptsize select one & \footnotesize While you are here for the season, where do you sleep?\par\vspace{1pt}{\footnotesize\color{hdr}{\hifont सीज़न में यहाँ रहते हुए आप कहाँ सोते हैं?}} & \scriptsize 1 At my own usual home (I live here) | 2 Rented room or house | 3 Room or dormitory the employer provides | 4 Sleeps at the shop, dhaba or workplace | 5 Tent or temporary shelter | 6 Lodge, dharamshala or ashram | 7 In the open, or a verandah | 8 Somewhere else & \scriptsize Lyons et al. 2023, Table 2 (area/settlement conditions), adapted to a labour-migrant setting\\
\footnotesize 92 & \footnotesize\texttt{electricity} & \scriptsize yes/no & \footnotesize Does your usual home have an electricity connection?\par\vspace{1pt}{\footnotesize\color{hdr}{\hifont क्या आपके असली घर में बिजली का कनेक्शन है?}} & \scriptsize 0 No | 1 Yes & \scriptsize NITI Aayog National MPI\\
\footnotesize 93 & \footnotesize\texttt{toilet\_type} & \scriptsize select one & \footnotesize What kind of toilet does your household use at your usual home?\par\vspace{1pt}{\footnotesize\color{hdr}{\hifont आपका घर अपने असली घर में किस तरह का शौचालय इस्तेमाल करता है?}} & \scriptsize 1 No toilet / open defecation | 2 Pit latrine without slab or open pit | 3 Improved toilet, but shared with other households | 4 Improved toilet, used only by this household & \scriptsize NITI Aayog National MPI sanitation indicator (improved/unimproved and shared/not)\\
\footnotesize 94 & \footnotesize\texttt{drinking\_water} & \scriptsize select one & \footnotesize What is your main source of drinking water at your usual home?\par\vspace{1pt}{\footnotesize\color{hdr}{\hifont आपके असली घर में पीने का पानी मुख्य रूप से कहाँ से आता है?}} & \scriptsize 1 Piped into the house or yard | 2 Public tap or standpipe | 3 Handpump, tubewell or borewell | 4 Protected well or protected spring | 5 Rainwater collection | 6 Bottled, packaged or community RO | 7 Unprotected well or unprotected spring | 8 River, stream, pond or canal | 9 Tanker truck or cart with drum & \scriptsize NITI Aayog National MPI drinking-water indicator (JMP improved/unimproved taxonomy)\\
\footnotesize 95 & \footnotesize\texttt{water\_on\_premises} & \scriptsize yes/no & \footnotesize Is the drinking water available at the house itself?\par\vspace{1pt}{\footnotesize\color{hdr}{\hifont क्या पीने का पानी घर पर ही मिल जाता है?}} & \scriptsize 0 No | 1 Yes & \scriptsize NITI Aayog National MPI drinking-water indicator (30-minute round-trip rule)\\
\footnotesize 96 & \footnotesize\texttt{water\_fetch\_minutes} & \scriptsize whole number & \footnotesize How long does it take to go there, get the water and come back? (minutes, round trip)\par\vspace{1pt}{\footnotesize\color{hdr}{\hifont वहाँ जाने, पानी भरने और लौटने में कुल कितना समय लगता है? (मिनट)}} & \scriptsize \textit{Asked only if: Ask only if water\_on\_premises = 0} & \scriptsize NITI Aayog National MPI drinking-water indicator (30-minute round-trip rule)\\
\footnotesize 97 & \footnotesize\texttt{cooking\_fuel} & \scriptsize select one & \footnotesize What does your household mainly use to cook at your usual home?\par\vspace{1pt}{\footnotesize\color{hdr}{\hifont आपके असली घर में खाना मुख्य रूप से किस चीज़ पर बनता है?}} & \scriptsize 1 LPG or cylinder gas | 2 Piped natural gas | 3 Electricity | 4 Biogas (gobar gas plant) | 5 Kerosene | 6 Firewood | 7 Dung cakes (gobar) | 8 Crop residue, straw or shrubs | 9 Charcoal | 10 Coal or lignite & \scriptsize NITI Aayog National MPI cooking-fuel indicator (its own list of dirty fuels)\\
\footnotesize 98 & \footnotesize\texttt{land\_unit} & \scriptsize select one & \footnotesize In what unit do you count your land — nali, bigha, acres or hectares?\par\vspace{1pt}{\footnotesize\color{hdr}{\hifont आप अपनी ज़मीन किस हिसाब से गिनते हैं — नाली, बीघा, एकड़ या हेक्टेयर?}} & \scriptsize 1 Nali | 2 Bigha | 3 Acres | 4 Hectares | 5 No land & \scriptsize Project design (local land units, Uttarakhand)\\
\footnotesize 99 & \footnotesize\texttt{land\_cultivable\_acres} & \scriptsize number & \footnotesize And how much CULTIVABLE land is that? Do not count the land the house stands on.\par\vspace{1pt}{\footnotesize\color{hdr}{\hifont आपके घर के पास कितनी जोतने लायक़ ज़मीन है, अपनी या जोतने की, एकड़ में? जिस ज़मीन पर घर बना है उसे मत गिनिए। (न हो तो 0)}} & \scriptsize  & \scriptsize Chaudhuri et al. 2002 asset covariate\\
\footnotesize 100 & \footnotesize\texttt{assets\_owned} & \scriptsize select multiple & \footnotesize Which of these does your household own? Read the list out and check every one they say: television, radio, bicycle, motorcycle or scooter, car/jeep/truck, telephone of any kind, computer or laptop, animal cart, refrigerator.\par\vspace{1pt}{\footnotesize\color{hdr}{\hifont इनमें से आपके घर में क्या-क्या है? सूची पढ़कर सुनाएँ और जो वे बताएँ, सब पर निशान लगाएँ: टीवी, रेडियो, साइकिल, मोटरसाइकिल या स्कूटर, कार/जीप/ट्रक, कोई भी फ़ोन, कंप्यूटर या लैपटॉप, जानवर से चलने वाली गाड़ी, फ़्रिज।}} & \scriptsize 1 Television | 2 Radio | 3 Bicycle | 4 Motorcycle or scooter | 5 Car, jeep or truck | 6 Telephone of any kind (mobile or landline) | 7 Computer or laptop | 8 Cart pulled by an animal | 9 Refrigerator\par\vspace{1pt}\textit{Asked only if: May be left blank if the household owns none of them} & \scriptsize NITI Aayog National MPI, Assets (weight 1/21)\\
\footnotesize 101 & \footnotesize\texttt{livestock\_owned} & \scriptsize select multiple & \footnotesize And which of these animals does your household own? Cows or buffaloes, goats or sheep, ponies or mules.\par\vspace{1pt}{\footnotesize\color{hdr}{\hifont और इनमें से कौन-कौन से जानवर आपके घर में हैं? गाय या भैंस, बकरी या भेड़, घोड़े या खच्चर।}} & \scriptsize 1 Cows or buffaloes | 2 Goats or sheep | 3 Ponies or mules\par\vspace{1pt}\textit{Asked only if: May be left blank if the household owns no animals} & \scriptsize Chaudhuri et al. 2002 asset covariate; pony/mule class is this project's exposure variable\\
\footnotesize 102 & \footnotesize\texttt{owns\_work\_vehicle} & \scriptsize yes/no & \footnotesize Do you use a vehicle to earn money — for example hauling goods, transporting people, or running a taxi? It does not have to be yours.\par\vspace{1pt}{\footnotesize\color{hdr}{\hifont क्या आपके घर की अपनी गाड़ी है जिससे आप कमाते हैं (जैसे जीप, ट्रक या टैक्सी)?}} & \scriptsize 0 No | 1 Yes & \scriptsize Chaudhuri et al. 2002 asset covariate\\
\footnotesize 103 & \footnotesize\texttt{owns\_work\_equipment} & \scriptsize yes/no & \footnotesize Does your household own equipment or tools that you use to earn?\par\vspace{1pt}{\footnotesize\color{hdr}{\hifont क्या आपके घर के पास काम के औज़ार या मशीन है जिससे आप कमाते हैं?}} & \scriptsize 0 No | 1 Yes & \scriptsize Chaudhuri et al. 2002 asset covariate\\
\footnotesize 104 & \footnotesize\texttt{work\_equipment\_detail} & \scriptsize free text & \footnotesize What are they? Write what they say. Leave blank if they cannot say.\par\vspace{1pt}{\footnotesize\color{hdr}{\hifont वे क्या हैं? जैसा वे बताएँ वैसा लिख लीजिए। न बता पाएँ तो खाली छोड़ दीजिए।}} & \scriptsize \textit{Asked only if: Ask only if owns\_work\_equipment = 1; may be left blank} & \scriptsize Standard household-survey item\\
\bottomrule\end{longtable}\vspace{6pt}
\noindent{\large\bfseries\color{hdr} Module G --- Finance, insurance and schemes}\quad{\small(16 questions)}\\[3pt]
{\small\itshape Read aloud before this module: Now about savings, loans, insurance and government schemes.}\\[5pt]
\begin{longtable}{@{}p{0.5cm}p{3.3cm}p{1.7cm}p{7.3cm}p{5.6cm}p{5.4cm}@{}}
\toprule \footnotesize\textbf{\#} & \footnotesize\textbf{Variable} & \footnotesize\textbf{Type} & \footnotesize\textbf{Question (English / Hindi)} & \footnotesize\textbf{Options {\&} skip} & \footnotesize\textbf{Source}\\ \midrule
\endhead
\footnotesize 105 & \footnotesize\texttt{has\_bank\_account} & \scriptsize yes/no & \footnotesize Does anyone in your household have a bank account or a post office account?\par\vspace{1pt}{\footnotesize\color{hdr}{\hifont क्या आपके घर में किसी का बैंक खाता या डाकघर का खाता है?}} & \scriptsize 0 No | 1 Yes & \scriptsize NITI Aayog National MPI bank-account indicator; VEP adaptive-capacity covariate\\
\footnotesize 106 & \footnotesize\texttt{took\_loan\_12m} & \scriptsize yes/no & \footnotesize In the last 12 months, did you or anyone in your household borrow money?\par\vspace{1pt}{\footnotesize\color{hdr}{\hifont पिछले 12 महीनों में क्या आपने या आपके घर के किसी सदस्य ने पैसा उधार लिया?}} & \scriptsize 0 No | 1 Yes & \scriptsize VEP adaptive-capacity covariate\\
\footnotesize 107 & \footnotesize\texttt{credit\_source} & \scriptsize select one & \footnotesize Who was the main lender?\par\vspace{1pt}{\footnotesize\color{hdr}{\hifont मुख्य रूप से किससे उधार लिया?}} & \scriptsize 1 Nationalised or public-sector bank | 2 Private bank | 3 Cooperative bank or RRB | 4 Microfinance institution or SHG | 5 Moneylender | 6 Relative or friend | 7 Employer or contractor (advance) | 8 Other\par\vspace{1pt}\textit{Asked only if: Ask only if took\_loan\_12m = 1} & \scriptsize VEP adaptive-capacity covariate (lender type)\\
\footnotesize 108 & \footnotesize\texttt{loan\_purpose} & \scriptsize select one & \footnotesize What did you take it for?\par\vspace{1pt}{\footnotesize\color{hdr}{\hifont वह क़र्ज़ आपने किस काम के लिए लिया था?}} & \scriptsize 1 Medical or hospital costs | 2 Buying an animal, vehicle or equipment for work | 3 Shop stock or business | 4 Food or daily household needs | 5 A wedding, funeral or ceremony | 6 House building or repair | 7 Education | 8 Repaying another loan | 9 Something else\par\vspace{1pt}\textit{Asked only if: Ask only if took\_loan\_12m = 1} & \scriptsize Project design (investment against distress borrowing)\\
\footnotesize 109 & \footnotesize\texttt{loan\_amount\_borrowed} & \scriptsize amount (Rs) & \footnotesize How much did you borrow in total? (Rs)\par\vspace{1pt}{\footnotesize\color{hdr}{\hifont आपने कुल कितना क़र्ज़ लिया था? (रुपये)}} & \scriptsize \textit{Asked only if: Ask only if took\_loan\_12m = 1} & \scriptsize Project design (loan size, for the debt-burden measures)\\
\footnotesize 110 & \footnotesize\texttt{loan\_amount} & \scriptsize amount (Rs) & \footnotesize How much of that loan is still left to repay? (Rs)\par\vspace{1pt}{\footnotesize\color{hdr}{\hifont उस क़र्ज़ में से अब भी कितना चुकाना बाक़ी है? (रुपये)}} & \scriptsize \textit{Asked only if: Ask only if took\_loan\_12m = 1} & \scriptsize Islam and Chowdhury 2025 (financial distress and household vulnerability to poverty)\\
\footnotesize 111 & \footnotesize\texttt{pays\_interest} & \scriptsize yes/no & \footnotesize Do you pay interest on it?\par\vspace{1pt}{\footnotesize\color{hdr}{\hifont क्या आप उस पर ब्याज देते हैं?}} & \scriptsize 0 No | 1 Yes\par\vspace{1pt}\textit{Asked only if: Ask only if took\_loan\_12m = 1} & \scriptsize Project design (informal lending is often interest-free)\\
\footnotesize 112 & \footnotesize\texttt{loan\_interest\_per100\_pm} & \scriptsize number & \footnotesize On every 100 rupees you borrowed, about how much interest do you pay in a month? (Rs; 0 if none, leave blank if not known)\par\vspace{1pt}{\footnotesize\color{hdr}{\hifont आपने जो 100 रुपये उधार लिए, उन पर महीने का कितना ब्याज देते हैं? (रुपये; दर पता न हो तो खाली छोड़ दें)}} & \scriptsize \textit{Asked only if: Ask only if took\_loan\_12m = 1; may be left blank} & \scriptsize Project design (informal-credit pricing as locally quoted)\\
\footnotesize 113 & \footnotesize\texttt{loan\_collateral} & \scriptsize select one & \footnotesize Did you have to give anything as security for it? If so, what — jewellery, land, animals, a vehicle, shop stock, a house, or something else?\par\vspace{1pt}{\footnotesize\color{hdr}{\hifont आपको क्या गिरवी रखना पड़ा?}} & \scriptsize 0 Nothing was given as security | 1 Jewellery or ornaments | 2 Land | 3 Animals | 4 A vehicle | 5 Shop stock or business goods | 6 House or building | 7 Something else\par\vspace{1pt}\textit{Asked only if: Ask only if took\_loan\_12m = 1} & \scriptsize Carter and Zimmerman 2000; Zimmerman and Carter 2003 (asset-based coping)\\
\footnotesize 114 & \footnotesize\texttt{n\_health\_insured} & \scriptsize whole number & \footnotesize How many people in your household are covered by any health insurance or health scheme, including government ones? (0 if none)\par\vspace{1pt}{\footnotesize\color{hdr}{\hifont आपके घर में कितने लोगों का कोई स्वास्थ्य बीमा या स्वास्थ्य योजना है, सरकारी भी मिलाकर? (कोई नहीं तो 0)}} & \scriptsize  & \scriptsize NITI Aayog National MPI health-insurance indicator; Lyons et al. 2023\\
\footnotesize 115 & \footnotesize\texttt{n\_life\_insured} & \scriptsize whole number & \footnotesize How many people in your household have life insurance? (0 if none)\par\vspace{1pt}{\footnotesize\color{hdr}{\hifont आपके घर में कितने लोगों का जीवन बीमा है? (कोई नहीं तो 0)}} & \scriptsize  & \scriptsize Standard household-survey item\\
\footnotesize 116 & \footnotesize\texttt{has\_crop\_insurance} & \scriptsize yes/no & \footnotesize Is your crop insured?\par\vspace{1pt}{\footnotesize\color{hdr}{\hifont क्या आपकी फ़सल का बीमा है?}} & \scriptsize 0 No | 1 Yes\par\vspace{1pt}\textit{Asked only if: Ask only if land\_cultivable\_acres > 0} & \scriptsize Standard household-survey item\\
\footnotesize 117 & \footnotesize\texttt{govt\_any\_benefit} & \scriptsize yes/no & \footnotesize In the last 12 months, did your household get anything from any government scheme? Read the list out before you take an answer.\par\vspace{1pt}{\footnotesize\color{hdr}{\hifont पिछले 12 महीनों में क्या आपके घर को किसी सरकारी योजना से कुछ मिला? जवाब लेने से पहले सूची पढ़कर सुनाइए।}} & \scriptsize 0 No | 1 Yes & \scriptsize Project design; scheme names as administered in Uttarakhand\\
\footnotesize 118 & \footnotesize\texttt{govt\_schemes\_detail} & \scriptsize free text & \footnotesize Which ones? Write down every scheme they name, in their own words.\par\vspace{1pt}{\footnotesize\color{hdr}{\hifont कौन-कौन सी? वे जितनी योजनाएँ बताएँ, सब उन्हीं के शब्दों में लिख लीजिए।}} & \scriptsize \textit{Asked only if: Ask only if govt\_any\_benefit = 1} & \scriptsize Project design; coded in the office against the Uttarakhand scheme list\\
\footnotesize 119 & \footnotesize\texttt{smartphone\_owned} & \scriptsize yes/no & \footnotesize Do you own a smartphone?\par\vspace{1pt}{\footnotesize\color{hdr}{\hifont क्या आपके पास स्मार्टफ़ोन है?}} & \scriptsize 0 No | 1 Yes & \scriptsize VEP adaptive-capacity covariate; STEP digital items\\
\footnotesize 120 & \footnotesize\texttt{n\_can\_transact\_online} & \scriptsize whole number & \footnotesize How many people in your household can make a payment by phone themselves, without help? (0 if none)\par\vspace{1pt}{\footnotesize\color{hdr}{\hifont आपके घर में कितने लोग बिना किसी की मदद के फ़ोन से पैसे भेज सकते हैं? (कोई नहीं तो 0)}} & \scriptsize  & \scriptsize Project design (intra-household financial capability)\\
\bottomrule\end{longtable}\vspace{6pt}
\noindent{\large\bfseries\color{hdr} Module H --- Health}\quad{\small(10 questions)}\\[3pt]
{\small\itshape Read aloud before this module: Now a few questions about health in your household. You may leave out any of these.}\\[5pt]
\begin{longtable}{@{}p{0.5cm}p{3.3cm}p{1.7cm}p{7.3cm}p{5.6cm}p{5.4cm}@{}}
\toprule \footnotesize\textbf{\#} & \footnotesize\textbf{Variable} & \footnotesize\textbf{Type} & \footnotesize\textbf{Question (English / Hindi)} & \footnotesize\textbf{Options {\&} skip} & \footnotesize\textbf{Source}\\ \midrule
\endhead
\footnotesize 121 & \footnotesize\texttt{morbidity\_15d} & \scriptsize yes/no & \footnotesize In the last 15 days, was anyone in your household ill?\par\vspace{1pt}{\footnotesize\color{hdr}{\hifont पिछले 15 दिनों में क्या आपके घर में कोई बीमार पड़ा?}} & \scriptsize 0 No | 1 Yes\par\vspace{1pt}\textit{Asked only if: May be left blank} & \scriptsize NSS health module (15-day recall) [source not held by this project; citation unverified as of 2026-09-28]\\
\footnotesize 122 & \footnotesize\texttt{morbidity\_coping\_15d} & \scriptsize select one & \footnotesize How did the household mainly cope with the cost of this: used savings; borrowed money; sold or pawned assets; cut other consumption; got help from relatives or community; or something else?\par\vspace{1pt}{\footnotesize\color{hdr}{\hifont इसका ख़र्च आपके घर ने मुख्य रूप से कैसे उठाया: जमा पैसे से; उधार लेकर; कुछ बेचकर या गिरवी रखकर; बाकी ख़र्च घटाकर; रिश्तेदारों या गाँव वालों की मदद से; या किसी और तरह?}} & \scriptsize 1 Used savings | 2 Borrowed money | 3 Sold or pawned assets | 4 Cut consumption | 5 Help from relatives/community | 6 Paid it out of normal earnings, nothing given up | 7 Something else (say what)\par\vspace{1pt}\textit{Asked only if: Ask only if morbidity\_15d = 1. May be left blank} & \scriptsize Project design (a reported illness with no cost/coping follow-up said nothing about its burden)\\
\footnotesize 123 & \footnotesize\texttt{morbidity\_cost\_15d} & \scriptsize amount (Rs) & \footnotesize About how much did the household spend on this in the last 15 days (medicine, doctor or clinic fees, travel for care)? (Rs)\par\vspace{1pt}{\footnotesize\color{hdr}{\hifont पिछले 15 दिनों में इस पर घर का लगभग कितना ख़र्च हुआ (दवाई, डॉक्टर या क्लीनिक की फ़ीस, आने-जाने का ख़र्च)? (रुपये)}} & \scriptsize \textit{Asked only if: Ask only if morbidity\_15d = 1. May be left blank} & \scriptsize Project design\\
\footnotesize 124 & \footnotesize\texttt{func\_limitation} & \scriptsize select one & \footnotesize Do you have difficulty walking or climbing steps?\par\vspace{1pt}{\footnotesize\color{hdr}{\hifont क्या आपको चलने या सीढ़ियाँ चढ़ने में दिक़्क़त होती है?}} & \scriptsize 1 No difficulty | 2 Some difficulty | 3 A lot of difficulty | 4 Cannot do it at all\par\vspace{1pt}\textit{Asked only if: May be left blank} & \scriptsize Washington Group on Disability Statistics, Short Set item 3 (mobility), verbatim; Lyons et al. 2023 health dimension\\
\footnotesize 125 & \footnotesize\texttt{hospitalization\_365d} & \scriptsize yes/no & \footnotesize In the last 12 months, was anyone in your household admitted to hospital overnight?\par\vspace{1pt}{\footnotesize\color{hdr}{\hifont पिछले 12 महीनों में क्या आपके घर का कोई रात भर के लिए अस्पताल में भर्ती हुआ?}} & \scriptsize 0 No | 1 Yes\par\vspace{1pt}\textit{Asked only if: May be left blank} & \scriptsize NSS health module (365-day recall) [source not held by this project; citation unverified as of 2026-09-28]\\
\footnotesize 126 & \footnotesize\texttt{health\_access\_barrier\_3m} & \scriptsize yes/no & \footnotesize In the last 3 months, was there a time when someone in your household needed medical care but could not get it?\par\vspace{1pt}{\footnotesize\color{hdr}{\hifont पिछले 3 महीनों में क्या कभी ऐसा हुआ कि आपके घर में किसी को इलाज की ज़रूरत थी पर मिल नहीं पाया?}} & \scriptsize 0 No | 1 Yes\par\vspace{1pt}\textit{Asked only if: May be left blank} & \scriptsize VASyR 2025 barriers\_health\_case\_access\_phc\_m (primary-care access barriers); Lyons et al. 2023 Table 2 indicator 2 'healthcare access'\\
\footnotesize 127 & \footnotesize\texttt{child\_death\_5y} & \scriptsize yes/no & \footnotesize In the last five years, has any child or young person under 18 in your household died?\par\vspace{1pt}{\footnotesize\color{hdr}{\hifont पिछले पाँच साल में क्या आपके घर के 18 साल से छोटे किसी बच्चे या किशोर की मृत्यु हुई?}} & \scriptsize 0 No | 1 Yes\par\vspace{1pt}\textit{Asked only if: May be left blank} & \scriptsize NITI Aayog National MPI, Child and Adolescent Mortality (weight 1/12)\\
\footnotesize 128 & \footnotesize\texttt{birth\_last\_5y} & \scriptsize yes/no & \footnotesize In the last five years, did any woman in your household give birth?\par\vspace{1pt}{\footnotesize\color{hdr}{\hifont पिछले पाँच साल में क्या आपके घर की किसी महिला ने बच्चे को जन्म दिया?}} & \scriptsize 0 No | 1 Yes\par\vspace{1pt}\textit{Asked only if: May be left blank} & \scriptsize NITI Aayog National MPI, Maternal Health (weight 1/12)\\
\footnotesize 129 & \footnotesize\texttt{anc\_4\_visits} & \scriptsize select one & \footnotesize For the most recent birth, did she have at least four check-ups before the delivery?\par\vspace{1pt}{\footnotesize\color{hdr}{\hifont सबसे हाल के जन्म के लिए, क्या डिलीवरी से पहले उनकी कम से कम चार बार जाँच हुई थी?}} & \scriptsize 0 No | 1 Yes\par\vspace{1pt}\textit{Asked only if: Ask only if birth\_last\_5y = 1. May be left blank} & \scriptsize NITI Aayog National MPI, Maternal Health (antenatal care limb)\\
\footnotesize 130 & \footnotesize\texttt{skilled\_birth\_attendant} & \scriptsize select one & \footnotesize Who conducted that delivery — a doctor, a nurse or ANM, an ASHA or Anganwadi worker, a dai, or nobody trained?\par\vspace{1pt}{\footnotesize\color{hdr}{\hifont वह डिलीवरी किसने कराई — डॉक्टर, नर्स या एएनएम, आशा या आँगनवाड़ी कार्यकर्ता, दाई, या कोई प्रशिक्षित व्यक्ति नहीं?}} & \scriptsize 1 Doctor | 2 Nurse, ANM or LHV | 3 ASHA or Anganwadi worker only | 4 Dai (traditional birth attendant) | 5 Nobody trained; family only\par\vspace{1pt}\textit{Asked only if: Ask only if birth\_last\_5y = 1. May be left blank} & \scriptsize NITI Aayog National MPI, Maternal Health (assisted delivery limb); skilled-provider definition per NFHS\\
\bottomrule\end{longtable}\vspace{6pt}
\noindent{\large\bfseries\color{hdr} Module I --- Food security, shocks and coping}\quad{\small(10 questions)}\\[3pt]
{\small\itshape Read aloud before this module: Now about food, and about any difficult times in the last year and how your household managed.}\\[5pt]
\begin{longtable}{@{}p{0.5cm}p{3.3cm}p{1.7cm}p{7.3cm}p{5.6cm}p{5.4cm}@{}}
\toprule \footnotesize\textbf{\#} & \footnotesize\textbf{Variable} & \footnotesize\textbf{Type} & \footnotesize\textbf{Question (English / Hindi)} & \footnotesize\textbf{Options {\&} skip} & \footnotesize\textbf{Source}\\ \midrule
\endhead
\footnotesize 131 & \footnotesize\texttt{distress\_event\_last365d} & \scriptsize select multiple & \footnotesize In the last 12 months, did any of these happen — to your household, or to the Yatra route you work on? Check every one that applies.\par\vspace{1pt}{\footnotesize\color{hdr}{\hifont पिछले 12 महीनों में इनमें से क्या-क्या हुआ — आपके घर के साथ, या जिस यात्रा मार्ग पर आप काम करते हैं उस पर? जो-जो हुआ, सब चुनिए।}} & \scriptsize 9 Nothing of this kind happened | 1 Serious illness or death of an earning member | 2 Crop failure or livestock loss | 3 Damage from a natural disaster (flood, landslide, fire) | 4 Loss of business, goods or assets | 5 Road or bridge blocked, route closed | 6 Yatra stopped early or badly disrupted | 7 Sharp fall in customers or prices | 8 Lost the work or the work stopped & \scriptsize VEP exposure covariate (Azeem et al. 2016); shock inventory following Gunther and Harttgen 2009, via Fujii 2016 section 4\\
\footnotesize 132 & \footnotesize\texttt{shock\_work\_lost\_weeks} & \scriptsize whole number & \footnotesize Thinking of whichever of those hit you hardest — about how many weeks of work did you lose because of it? (0 if none)\par\vspace{1pt}{\footnotesize\color{hdr}{\hifont इनमें से जो सबसे भारी पड़ा — उसके कारण लगभग कितने हफ़्ते का काम छूटा? (कुछ न छूटा हो तो 0)}} & \scriptsize \textit{Asked only if: Ask only if distress\_event\_last365d names a real shock (not code 9, nothing happened)} & \scriptsize Ligon and Schechter 2003; Dercon and Krishnan 2000 (VER needs shock magnitude, not only incidence)\\
\footnotesize 133 & \footnotesize\texttt{shock\_money\_spent} & \scriptsize amount (Rs) & \footnotesize And about how much money did your household have to spend because of it — treatment, repairs, replacing what was lost? (Rs, 0 if nothing)\par\vspace{1pt}{\footnotesize\color{hdr}{\hifont और उसके कारण आपके घर को लगभग कितना पैसा ख़र्च करना पड़ा — इलाज, मरम्मत, जो नुक़सान हुआ उसे पूरा करना? (रुपये, कुछ नहीं तो 0)}} & \scriptsize \textit{Asked only if: Ask only if distress\_event\_last365d names a real shock (not code 9, nothing happened)} & \scriptsize Ligon and Schechter 2003; Dercon and Krishnan 2000 (VER needs shock magnitude, not only incidence)\\
\footnotesize 134 & \footnotesize\texttt{shock\_month} & \scriptsize select one & \footnotesize And in which month did that one happen?\par\vspace{1pt}{\footnotesize\color{hdr}{\hifont और वह किस महीने हुआ था?}} & \scriptsize 1 January | 2 February | 3 March | 4 April | 5 May | 6 June | 7 July | 8 August | 9 September | 10 October | 11 November | 12 December\par\vspace{1pt}\textit{Asked only if: Ask only if distress\_event\_last365d names a real shock (not code 9, nothing happened)} & \scriptsize Dercon and Krishnan 2000 (seasonal timing of shocks); project design\\
\footnotesize 135 & \footnotesize\texttt{shock\_coping} & \scriptsize select multiple & \footnotesize How did your household manage? Check every one they used.\par\vspace{1pt}{\footnotesize\color{hdr}{\hifont आपके घर ने कैसे सँभाला? उन्होंने जो-जो किया, सब चुनिए।}} & \scriptsize 1 Used savings | 2 Borrowed money | 3 Sold or pawned assets | 4 Cut consumption | 5 Help from relatives/community | 6 Paid it out of normal earnings, nothing given up | 7 Something else (say what)\par\vspace{1pt}\textit{Asked only if: Ask only if any shock was reported} & \scriptsize Coping strategies in VEP studies; Carter and Zimmerman 2000 on asset versus consumption smoothing\\
\footnotesize 136 & \footnotesize\texttt{fies\_skipped\_meal} & \scriptsize yes/no & \footnotesize In the last 12 months, was there a time when, because of a lack of money or other resources, you had to skip a meal?\par\vspace{1pt}{\footnotesize\color{hdr}{\hifont पिछले 12 महीनों में, पैसे या साधन की कमी के कारण, कभी ऐसा हुआ कि आपको एक वक़्त का खाना छोड़ना पड़ा?}} & \scriptsize 0 No | 1 Yes\par\vspace{1pt}\textit{Asked only if: May be left blank} & \scriptsize FAO Food Insecurity Experience Scale (FIES), SDG indicator 2.1.2; behavioural items only\\
\footnotesize 137 & \footnotesize\texttt{fies\_ate\_less} & \scriptsize yes/no & \footnotesize In the last 12 months, was there a time when, because of a lack of money or other resources, you ate less than you thought you should?\par\vspace{1pt}{\footnotesize\color{hdr}{\hifont पिछले 12 महीनों में, पैसे या साधन की कमी के कारण, कभी ऐसा हुआ कि आपने जितना खाना चाहिए था उससे कम खाया?}} & \scriptsize 0 No | 1 Yes\par\vspace{1pt}\textit{Asked only if: May be left blank} & \scriptsize FAO Food Insecurity Experience Scale (FIES), SDG indicator 2.1.2; behavioural items only\\
\footnotesize 138 & \footnotesize\texttt{fies\_ran\_out} & \scriptsize yes/no & \footnotesize In the last 12 months, was there a time when, because of a lack of money or other resources, your household ran out of food?\par\vspace{1pt}{\footnotesize\color{hdr}{\hifont पिछले 12 महीनों में, पैसे या साधन की कमी के कारण, कभी ऐसा हुआ कि आपके घर में खाना ख़त्म हो गया?}} & \scriptsize 0 No | 1 Yes\par\vspace{1pt}\textit{Asked only if: May be left blank} & \scriptsize FAO Food Insecurity Experience Scale (FIES), SDG indicator 2.1.2; behavioural items only\\
\footnotesize 139 & \footnotesize\texttt{fies\_hungry} & \scriptsize yes/no & \footnotesize In the last 12 months, was there a time when, because of a lack of money or other resources, you were hungry but did not eat?\par\vspace{1pt}{\footnotesize\color{hdr}{\hifont पिछले 12 महीनों में, पैसे या साधन की कमी के कारण, कभी ऐसा हुआ कि आपको भूख लगी पर आपने खाया नहीं?}} & \scriptsize 0 No | 1 Yes\par\vspace{1pt}\textit{Asked only if: May be left blank} & \scriptsize FAO Food Insecurity Experience Scale (FIES), SDG indicator 2.1.2; behavioural items only\\
\footnotesize 140 & \footnotesize\texttt{fies\_whole\_day} & \scriptsize yes/no & \footnotesize In the last 12 months, was there a time when, because of a lack of money or other resources, you went without eating for a whole day?\par\vspace{1pt}{\footnotesize\color{hdr}{\hifont पिछले 12 महीनों में, पैसे या साधन की कमी के कारण, कभी ऐसा हुआ कि आपने पूरा दिन कुछ नहीं खाया?}} & \scriptsize 0 No | 1 Yes\par\vspace{1pt}\textit{Asked only if: May be left blank} & \scriptsize FAO Food Insecurity Experience Scale (FIES), SDG indicator 2.1.2; behavioural items only\\
\bottomrule\end{longtable}\vspace{6pt}
\noindent{\large\bfseries\color{hdr} Module J --- Ropeway}\quad{\small(5 questions)}\\[3pt]
{\small\itshape Read aloud before this module: Now a few questions about the ropeway that has been proposed for this route. It is only a proposal — nothing has been built. There is no right answer and nothing you say here goes to anyone but us.}\\[5pt]
\begin{longtable}{@{}p{0.5cm}p{3.3cm}p{1.7cm}p{7.3cm}p{5.6cm}p{5.4cm}@{}}
\toprule \footnotesize\textbf{\#} & \footnotesize\textbf{Variable} & \footnotesize\textbf{Type} & \footnotesize\textbf{Question (English / Hindi)} & \footnotesize\textbf{Options {\&} skip} & \footnotesize\textbf{Source}\\ \midrule
\endhead
\footnotesize 141 & \footnotesize\texttt{ropeway\_stance} & \scriptsize select one & \footnotesize A ropeway is proposed for this route. Are you in favour, neutral, or against?\par\vspace{1pt}{\footnotesize\color{hdr}{\hifont गौरीकुंड से केदारनाथ तक रोपवे बनाने की बात चल रही है। आप इसके पक्ष में हैं, बीच में हैं, या ख़िलाफ़ हैं?}} & \scriptsize 1 Support | 2 Neutral | 3 Oppose | 98 Prefer not to say & \scriptsize Project design\\
\footnotesize 142 & \footnotesize\texttt{ropeway\_expect\_work} & \scriptsize select one & \footnotesize If the ropeway is built, what do you think will happen to your own work? More work for you; less work; about the same; or don't know?\par\vspace{1pt}{\footnotesize\color{hdr}{\hifont अगर रोपवे बन गया, तो आपके अपने काम पर क्या असर होगा? मेरे लिए ज़्यादा काम; मेरे लिए कम काम; लगभग वैसा ही; या पता नहीं?}} & \scriptsize 1 More work for me | 2 Less work for me | 3 About the same | 97 Don't know & \scriptsize Pelz et al. 2024, Energy Policy 186:113973 (concept: perceived effect on livelihoods; wording ours)\\
\footnotesize 143 & \footnotesize\texttt{ropeway\_expect\_jobs} & \scriptsize select one & \footnotesize Who do you think will get most of the jobs the ropeway creates? People from these villages; people from outside the area; nobody will get jobs; or don't know?\par\vspace{1pt}{\footnotesize\color{hdr}{\hifont रोपवे में जो नौकरियाँ बनेंगी, वे ज़्यादातर किसे मिलेंगी? इन गाँवों के लोगों को; इलाके के बाहर के लोगों को; किसी को नहीं मिलेंगी; या पता नहीं?}} & \scriptsize 1 People from these villages | 2 People from outside the area | 3 Nobody will get jobs | 97 Don't know & \scriptsize Pelz et al. 2024, Energy Policy 186:113973 (concept: local vs outside benefit; wording ours)\\
\footnotesize 144 & \footnotesize\texttt{ropeway\_trust} & \scriptsize select one & \footnotesize Do you think the ropeway will improve the livelihoods of people here? Yes; no; or don't know?\par\vspace{1pt}{\footnotesize\color{hdr}{\hifont क्या आपको लगता है कि रोपवे से यहाँ के लोगों की रोज़ी-रोटी सुधरेगी? हाँ; नहीं; या पता नहीं?}} & \scriptsize 0 No | 1 Yes | 97 Don't know & \scriptsize Pelz et al. 2024, Energy Policy 186:113973 (Fig. 4B, trust in the company; wording ours)\\
\footnotesize 145 & \footnotesize\texttt{trek\_dependent} & \scriptsize yes/no & \footnotesize Is your main work carrying, transporting or guiding people or goods on foot along the pilgrimage trek?\par\vspace{1pt}{\footnotesize\color{hdr}{\hifont क्या आपका मुख्य यात्रा-काम गौरीकुंड से केदारनाथ के पैदल रास्ते पर लोगों या सामान को ले जाना, पहुँचाना या रास्ता दिखाना है?}} & \scriptsize 0 No | 1 Yes & \scriptsize Project design (see tasks\_module W4: exposure must be asked directly)\\
\bottomrule\end{longtable}\vspace{6pt}
\noindent{\large\bfseries\color{hdr} Module K --- Job quality (Apablaza et al. 2026 questions)}\quad{\small(22 questions)}\\[3pt]
{\small\itshape Read aloud before this module: Now a set of questions about the conditions of your work — hours, pay, contract and so on. These come from a standard list used in many countries, so one or two may not fit your situation. Say so and we will move on.}\\[5pt]
\begin{longtable}{@{}p{0.5cm}p{3.3cm}p{1.7cm}p{7.3cm}p{5.6cm}p{5.4cm}@{}}
\toprule \footnotesize\textbf{\#} & \footnotesize\textbf{Variable} & \footnotesize\textbf{Type} & \footnotesize\textbf{Question (English / Hindi)} & \footnotesize\textbf{Options {\&} skip} & \footnotesize\textbf{Source}\\ \midrule
\endhead
\footnotesize 146 & \footnotesize\texttt{contract\_status} & \scriptsize select one & \footnotesize Do you have a signed contract? Yes, signed; yes but not yet signed; no contract.\par\vspace{1pt}{\footnotesize\color{hdr}{\hifont क्या आपका दस्तख़त किया हुआ अनुबंध (कॉन्ट्रैक्ट) है? हाँ, दस्तख़त किया; हाँ, पर अभी दस्तख़त नहीं हुआ; कोई अनुबंध नहीं।}} & \scriptsize 1 Yes, signed | 2 Yes, but not yet signed | 3 No contract\par\vspace{1pt}\textit{Asked only if: Ask only if employment\_type is 3 or 4 (wage workers) — self-employed skip to the next question} & \scriptsize Apablaza et al. 2026, Appendix 2 Q11\\
\footnotesize 147 & \footnotesize\texttt{workplace\_registered} & \scriptsize yes/no & \footnotesize Is your workplace or business registered, for example with a taxpayer or GST number, a shop or trade licence, the Yatra registration, or a union?\par\vspace{1pt}{\footnotesize\color{hdr}{\hifont क्या आपका काम या धंधा रजिस्टर्ड है, जैसे GST या टैक्स नंबर, दुकान का लाइसेंस, यात्रा का रजिस्ट्रेशन, या कोई यूनियन?}} & \scriptsize 0 No | 1 Yes & \scriptsize Apablaza et al. 2026, Appendix 2 Q12 (widened to local registrations)\\
\footnotesize 148 & \footnotesize\texttt{pension\_contrib} & \scriptsize select one & \footnotesize Do you contribute to any pension system? Yes, the employer deducts it; yes, voluntarily; no.\par\vspace{1pt}{\footnotesize\color{hdr}{\hifont क्या आप किसी पेंशन में पैसा डालते हैं? हाँ, मालिक तनख़्वाह से काटता है; हाँ, अपनी मर्ज़ी से डालते हैं; नहीं।}} & \scriptsize 1 Yes, employer deducts it | 2 Yes, contributes voluntarily | 3 No & \scriptsize Apablaza et al. 2026, Appendix 2 Q16\\
\footnotesize 149 & \footnotesize\texttt{work\_health\_ins} & \scriptsize select one & \footnotesize Do you have health insurance through your work? Yes; only private or other insurance; none; don't know.\par\vspace{1pt}{\footnotesize\color{hdr}{\hifont क्या काम की तरफ़ से आपका स्वास्थ्य बीमा है? हाँ; सिर्फ़ अपना या कोई और बीमा; कोई नहीं; पता नहीं।}} & \scriptsize 1 Yes, through work | 2 Only private or other insurance | 3 None | 4 Don't know & \scriptsize Apablaza et al. 2026, Appendix 2 Q17\\
\footnotesize 150 & \footnotesize\texttt{leave\_rights} & \scriptsize select one & \footnotesize Do you have the right to paid holiday, sick or maternity leave?\par\vspace{1pt}{\footnotesize\color{hdr}{\hifont क्या आपको पैसे वाली छुट्टी मिलती है — त्योहार, बीमारी या ज़च्चगी की?}} & \scriptsize 0 No | 1 Yes | 97 Don't know & \scriptsize Apablaza et al. 2026, Appendix 2 Q18\\
\footnotesize 151 & \footnotesize\texttt{injured\_ever} & \scriptsize select one & \footnotesize Have you ever been physically injured at your workplace?\par\vspace{1pt}{\footnotesize\color{hdr}{\hifont क्या आपको कभी काम की जगह पर चोट लगी है?}} & \scriptsize 0 No | 1 Yes | 97 Don't know & \scriptsize Apablaza et al. 2026, Appendix 2 Q19\\
\footnotesize 152 & \footnotesize\texttt{workplace\_injury\_12m} & \scriptsize select one & \footnotesize In the last 12 months, was anyone physically injured at your workplace because of work?\par\vspace{1pt}{\footnotesize\color{hdr}{\hifont पिछले 12 महीनों में क्या आपके काम की जगह पर काम की वजह से किसी को चोट लगी?}} & \scriptsize 0 No | 1 Yes | 97 Don't know & \scriptsize Apablaza et al. 2026, Appendix 2 Q20 (split into injury and death)\\
\footnotesize 153 & \footnotesize\texttt{workplace\_death\_12m} & \scriptsize select one & \footnotesize And in the last 12 months, did anyone die at your workplace because of work?\par\vspace{1pt}{\footnotesize\color{hdr}{\hifont और पिछले 12 महीनों में, क्या आपके काम की जगह पर काम की वजह से किसी की मौत हुई?}} & \scriptsize 0 No | 1 Yes | 97 Don't know & \scriptsize Apablaza et al. 2026, Appendix 2 Q20 (split into injury and death)\\
\footnotesize 154 & \footnotesize\texttt{wants\_more\_work} & \scriptsize select one & \footnotesize Would you like to work more hours than you do?\par\vspace{1pt}{\footnotesize\color{hdr}{\hifont क्या आप अभी से ज़्यादा घंटे काम करना चाहेंगे?}} & \scriptsize 0 No | 1 Yes | 97 Don't know & \scriptsize Apablaza et al. 2026, Appendix 2 Q21\\
\footnotesize 155 & \footnotesize\texttt{more\_hours\_day} & \scriptsize whole number & \footnotesize On a working day, how many MORE hours would you like to work?\par\vspace{1pt}{\footnotesize\color{hdr}{\hifont काम के दिन में आप कितने घंटे और काम करना चाहेंगे?}} & \scriptsize \textit{Asked only if: Ask only if wants\_more\_work = 1} & \scriptsize Apablaza et al. 2026, Appendix 2 Q22\\
\footnotesize 156 & \footnotesize\texttt{months\_looked\_for\_work} & \scriptsize whole number & \footnotesize Across the months when you had no paid work, about how many MONTHS in total were you looking for work?\par\vspace{1pt}{\footnotesize\color{hdr}{\hifont जिन महीनों में आपके पास कोई पैसे वाला काम नहीं था, उनमें कुल मिलाकर लगभग कितने महीने आप काम ढूँढते रहे?}} & \scriptsize \textit{Asked only if: Ask only if months\_no\_paid\_work > 0} & \scriptsize Apablaza et al. 2026, Appendix 2 Q23 (duration in weeks, tied to the calendar instead of a single point in time)\\
\footnotesize 157 & \footnotesize\texttt{first\_job\_ever} & \scriptsize select one & \footnotesize Was your current work the first paid job you ever had?\par\vspace{1pt}{\footnotesize\color{hdr}{\hifont क्या यह आपका पहला पैसे वाला काम है?}} & \scriptsize 0 No | 1 Yes | 97 Don't know & \scriptsize Apablaza et al. 2026, Appendix 2 Q24\\
\footnotesize 158 & \footnotesize\texttt{higher\_ed} & \scriptsize select one & \footnotesize Have you completed any college, diploma or university course?\par\vspace{1pt}{\footnotesize\color{hdr}{\hifont क्या आपने कॉलेज, डिप्लोमा या यूनिवर्सिटी का कोई कोर्स पूरा किया है?}} & \scriptsize 0 No | 1 Yes | 97 Don't know & \scriptsize Apablaza et al. 2026, Appendix 2 employment table, occupational status (1/8)\\
\footnotesize 159 & \footnotesize\texttt{social\_security\_any} & \scriptsize select one & \footnotesize Are you enrolled in any pension, insurance or social-security scheme? For example, EPFO or PM-SYM pension, or Ayushman Bharat health cover.\par\vspace{1pt}{\footnotesize\color{hdr}{\hifont क्या आप किसी पेंशन, बीमा या सरकारी सामाजिक सुरक्षा योजना में शामिल हैं? जैसे EPFO या PM-SYM पेंशन, या आयुष्मान भारत स्वास्थ्य कवर।}} & \scriptsize 0 No | 1 Yes | 97 Don't know & \scriptsize Apablaza et al. 2026, Appendix 2 employment table, employment security (1/8)\\
\footnotesize 160 & \footnotesize\texttt{intens\_speed} & \scriptsize select one & \footnotesize For more than half of your working day, do you work at very high speed?\par\vspace{1pt}{\footnotesize\color{hdr}{\hifont काम के दिन के आधे से ज़्यादा समय, क्या आप बहुत तेज़ रफ़्तार से काम करते हैं?}} & \scriptsize 0 No | 1 Yes | 97 Don't know & \scriptsize Apablaza et al. 2026, Appendix 2 employment table, work intensity (1/16)\\
\footnotesize 161 & \footnotesize\texttt{intens\_deadline} & \scriptsize select one & \footnotesize For more than half of your working day, do you work to tight deadlines?\par\vspace{1pt}{\footnotesize\color{hdr}{\hifont काम के दिन के आधे से ज़्यादा समय, क्या आपको कड़ी समय-सीमा में काम करना पड़ता है?}} & \scriptsize 0 No | 1 Yes | 97 Don't know & \scriptsize Apablaza et al. 2026, Appendix 2 employment table, work intensity (1/16)\\
\footnotesize 162 & \footnotesize\texttt{intens\_time} & \scriptsize select one & \footnotesize For more than half of your working day, do you not have enough time to finish your tasks?\par\vspace{1pt}{\footnotesize\color{hdr}{\hifont काम के दिन के आधे से ज़्यादा समय, क्या आपके पास काम पूरा करने के लिए काफ़ी समय नहीं होता?}} & \scriptsize 0 No | 1 Yes | 97 Don't know & \scriptsize Apablaza et al. 2026, Appendix 2 employment table, work intensity (1/16)\\
\footnotesize 163 & \footnotesize\texttt{posture\_position} & \scriptsize select one & \footnotesize For more than half of your working day, do you work in a tiring or painful position?\par\vspace{1pt}{\footnotesize\color{hdr}{\hifont काम के दिन के आधे से ज़्यादा समय, क्या आप थकाने वाली या दर्द वाली स्थिति में काम करते हैं?}} & \scriptsize 0 No | 1 Yes | 97 Don't know & \scriptsize Apablaza et al. 2026, Appendix 2 employment table, posture risk (1/16)\\
\footnotesize 164 & \footnotesize\texttt{posture\_loads} & \scriptsize select one & \footnotesize For more than half of your working day, do you carry or move heavy loads?\par\vspace{1pt}{\footnotesize\color{hdr}{\hifont काम के दिन के आधे से ज़्यादा समय, क्या आप भारी बोझ उठाते या ढोते हैं?}} & \scriptsize 0 No | 1 Yes | 97 Don't know & \scriptsize Apablaza et al. 2026, Appendix 2 employment table, posture risk (1/16)\\
\footnotesize 165 & \footnotesize\texttt{posture\_repetitive} & \scriptsize select one & \footnotesize For more than half of your working day, do you make the same movements again and again?\par\vspace{1pt}{\footnotesize\color{hdr}{\hifont काम के दिन के आधे से ज़्यादा समय, क्या आप एक ही तरह की हरकत बार-बार करते हैं?}} & \scriptsize 0 No | 1 Yes | 97 Don't know & \scriptsize Apablaza et al. 2026, Appendix 2 employment table, posture risk (1/16)\\
\footnotesize 166 & \footnotesize\texttt{phys\_noise} & \scriptsize select one & \footnotesize For more than half of your working day, are you exposed to loud noise?\par\vspace{1pt}{\footnotesize\color{hdr}{\hifont काम के दिन के आधे से ज़्यादा समय, क्या आप तेज़ शोर में रहते हैं?}} & \scriptsize 0 No | 1 Yes | 97 Don't know & \scriptsize Apablaza et al. 2026, Appendix 2 employment table, physical risk (1/16)\\
\footnotesize 167 & \footnotesize\texttt{phys\_temperature} & \scriptsize select one & \footnotesize For more than half of your working day, are you exposed to extreme heat or cold?\par\vspace{1pt}{\footnotesize\color{hdr}{\hifont काम के दिन के आधे से ज़्यादा समय, क्या आप बहुत गर्मी या बहुत ठंड में रहते हैं?}} & \scriptsize 0 No | 1 Yes | 97 Don't know & \scriptsize Apablaza et al. 2026, Appendix 2 employment table, physical risk (1/16)\\
\bottomrule\end{longtable}\vspace{6pt}
\noindent{\large\bfseries\color{hdr} Module L --- Tasks and skills (short block)}\quad{\small(17 questions)}\\[3pt]
{\small\itshape Read aloud before this module: Now I will read out some kinds of work-tasks and ask whether you do them. There is no right or wrong answer — we are trying to understand what skills the work here actually uses.}\\[5pt]
\begin{longtable}{@{}p{0.5cm}p{3.3cm}p{1.7cm}p{7.3cm}p{5.6cm}p{5.4cm}@{}}
\toprule \footnotesize\textbf{\#} & \footnotesize\textbf{Variable} & \footnotesize\textbf{Type} & \footnotesize\textbf{Question (English / Hindi)} & \footnotesize\textbf{Options {\&} skip} & \footnotesize\textbf{Source}\\ \midrule
\endhead
\footnotesize 168 & \footnotesize\texttt{tk\_load} & \scriptsize select one & \footnotesize Do you load, unload, stack or count goods? \par\vspace{1pt}{\footnotesize\color{hdr}{\hifont क्या आप सामान चढ़ाना-उतारना, लगाना या गिनना का काम करते हैं? (1 अपने मुख्य यात्रा-काम में नियमित; 2 मुख्य काम में नहीं, पर पहले कहीं और किया है; 3 कभी नहीं)}} & \scriptsize 1 Yes, regularly in my main Yatra work | 2 Not in my main work, but done before elsewhere | 3 Never done & \scriptsize Gathmann and Schonberg 2010; NCO 9333.0100 loader\\
\footnotesize 169 & \footnotesize\texttt{tk\_drive} & \scriptsize select one & \footnotesize Do you drive a motor vehicle? \par\vspace{1pt}{\footnotesize\color{hdr}{\hifont क्या आप गाड़ी चलाना का काम करते हैं? (1 अपने मुख्य यात्रा-काम में नियमित; 2 मुख्य काम में नहीं, पर पहले कहीं और किया है; 3 कभी नहीं)}} & \scriptsize 1 Yes, regularly in my main Yatra work | 2 Not in my main work, but done before elsewhere | 3 Never done & \scriptsize Spitz-Oener 2006 / STEP; NCO 8322.0100\\
\footnotesize 170 & \footnotesize\texttt{tk\_engine} & \scriptsize select one & \footnotesize Do you operate an engine or machine? \par\vspace{1pt}{\footnotesize\color{hdr}{\hifont क्या आप इंजन या मशीन चलाना का काम करते हैं? (1 अपने मुख्य यात्रा-काम में नियमित; 2 मुख्य काम में नहीं, पर पहले कहीं और किया है; 3 कभी नहीं)}} & \scriptsize 1 Yes, regularly in my main Yatra work | 2 Not in my main work, but done before elsewhere | 3 Never done & \scriptsize Gathmann and Schonberg 2010; NCO 8343.1700 ropeway operator\\
\footnotesize 171 & \footnotesize\texttt{tk\_electric} & \scriptsize select one & \footnotesize Do you do electrical or wiring work? \par\vspace{1pt}{\footnotesize\color{hdr}{\hifont क्या आप बिजली या तार का काम करना का काम करते हैं? (1 अपने मुख्य यात्रा-काम में नियमित; 2 मुख्य काम में नहीं, पर पहले कहीं और किया है; 3 कभी नहीं)}} & \scriptsize 1 Yes, regularly in my main Yatra work | 2 Not in my main work, but done before elsewhere | 3 Never done & \scriptsize Spitz-Oener 2006; NCO 7411.0301 wireman\\
\footnotesize 172 & \footnotesize\texttt{tk\_safety} & \scriptsize select one & \footnotesize Do you check equipment or the route for safety? \par\vspace{1pt}{\footnotesize\color{hdr}{\hifont क्या आप सामान या रास्ते की सुरक्षा जाँचना का काम करते हैं? (1 अपने मुख्य यात्रा-काम में नियमित; 2 मुख्य काम में नहीं, पर पहले कहीं और किया है; 3 कभी नहीं)}} & \scriptsize 1 Yes, regularly in my main Yatra work | 2 Not in my main work, but done before elsewhere | 3 Never done & \scriptsize Gathmann and Schonberg 2010; NCO 8343.1700\\
\footnotesize 173 & \footnotesize\texttt{tk\_sell} & \scriptsize select one & \footnotesize Do you sell goods or services? \par\vspace{1pt}{\footnotesize\color{hdr}{\hifont क्या आप सामान या सेवा बेचना का काम करते हैं? (1 अपने मुख्य यात्रा-काम में नियमित; 2 मुख्य काम में नहीं, पर पहले कहीं और किया है; 3 कभी नहीं)}} & \scriptsize 1 Yes, regularly in my main Yatra work | 2 Not in my main work, but done before elsewhere | 3 Never done & \scriptsize Gathmann and Schonberg 2010; NCO 5223.0200\\
\footnotesize 174 & \footnotesize\texttt{tk\_cash} & \scriptsize select one & \footnotesize Do you handle cash and payments? \par\vspace{1pt}{\footnotesize\color{hdr}{\hifont क्या आप पैसे और भुगतान सँभालना का काम करते हैं? (1 अपने मुख्य यात्रा-काम में नियमित; 2 मुख्य काम में नहीं, पर पहले कहीं और किया है; 3 कभी नहीं)}} & \scriptsize 1 Yes, regularly in my main Yatra work | 2 Not in my main work, but done before elsewhere | 3 Never done & \scriptsize Autor and Handel 2013; NCO 5131.0401\\
\footnotesize 175 & \footnotesize\texttt{tk\_cook} & \scriptsize select one & \footnotesize Do you cook or prepare food or drink? \par\vspace{1pt}{\footnotesize\color{hdr}{\hifont क्या आप खाना या पीने की चीज़ बनाना का काम करते हैं? (1 अपने मुख्य यात्रा-काम में नियमित; 2 मुख्य काम में नहीं, पर पहले कहीं और किया है; 3 कभी नहीं)}} & \scriptsize 1 Yes, regularly in my main Yatra work | 2 Not in my main work, but done before elsewhere | 3 Never done & \scriptsize NCO 5120.0300 cook\\
\footnotesize 176 & \footnotesize\texttt{tk\_serve} & \scriptsize select one & \footnotesize Do you serve guests or customers? \par\vspace{1pt}{\footnotesize\color{hdr}{\hifont क्या आप मेहमानों या ग्राहकों की सेवा करना का काम करते हैं? (1 अपने मुख्य यात्रा-काम में नियमित; 2 मुख्य काम में नहीं, पर पहले कहीं और किया है; 3 कभी नहीं)}} & \scriptsize 1 Yes, regularly in my main Yatra work | 2 Not in my main work, but done before elsewhere | 3 Never done & \scriptsize NCO 5131.0401 waiter; Pal et al. 2026\\
\footnotesize 177 & \footnotesize\texttt{tk\_clean} & \scriptsize select one & \footnotesize Do you clean rooms or public areas? \par\vspace{1pt}{\footnotesize\color{hdr}{\hifont क्या आप कमरे या सार्वजनिक जगह साफ़ करना का काम करते हैं? (1 अपने मुख्य यात्रा-काम में नियमित; 2 मुख्य काम में नहीं, पर पहले कहीं और किया है; 3 कभी नहीं)}} & \scriptsize 1 Yes, regularly in my main Yatra work | 2 Not in my main work, but done before elsewhere | 3 Never done & \scriptsize NCO 5151.0202 room attendant; Autor-Levy-Murnane 2003 (janitorial)\\
\footnotesize 178 & \footnotesize\texttt{tk\_guide} & \scriptsize select one & \footnotesize Do you guide or explain things to visitors? \par\vspace{1pt}{\footnotesize\color{hdr}{\hifont क्या आप आने वालों को रास्ता या बात समझाना का काम करते हैं? (1 अपने मुख्य यात्रा-काम में नियमित; 2 मुख्य काम में नहीं, पर पहले कहीं और किया है; 3 कभी नहीं)}} & \scriptsize 1 Yes, regularly in my main Yatra work | 2 Not in my main work, but done before elsewhere | 3 Never done & \scriptsize Gathmann and Schonberg 2010; NCO 5113.0200\\
\footnotesize 179 & \footnotesize\texttt{tk\_coord} & \scriptsize select one & \footnotesize Do you coordinate work by phone, radio or signals? \par\vspace{1pt}{\footnotesize\color{hdr}{\hifont क्या आप फ़ोन, रेडियो या इशारों से दूसरों के साथ तालमेल करना का काम करते हैं? (1 अपने मुख्य यात्रा-काम में नियमित; 2 मुख्य काम में नहीं, पर पहले कहीं और किया है; 3 कभी नहीं)}} & \scriptsize 1 Yes, regularly in my main Yatra work | 2 Not in my main work, but done before elsewhere | 3 Never done & \scriptsize NCO 8343.1700 ropeway operator\\
\footnotesize 180 & \footnotesize\texttt{drove\_twowheeler} & \scriptsize yes/no & \footnotesize Have you ever driven a motorcycle or scooter, with or without a licence?\par\vspace{1pt}{\footnotesize\color{hdr}{\hifont क्या आपने कभी मोटरसाइकिल या स्कूटर चलाई है, लाइसेंस हो या न हो?}} & \scriptsize 0 No | 1 Yes\par\vspace{1pt}\textit{Asked only if: Ask only if tk\_drive is 1 or 2} & \scriptsize Level ladder without licences: people drive without holding a licence\\
\footnotesize 181 & \footnotesize\texttt{drove\_car} & \scriptsize yes/no & \footnotesize ... a car or jeep?\par\vspace{1pt}{\footnotesize\color{hdr}{\hifont ...कार या जीप?}} & \scriptsize 0 No | 1 Yes\par\vspace{1pt}\textit{Asked only if: Ask only if tk\_drive is 1 or 2} & \scriptsize Level ladder without licences\\
\footnotesize 182 & \footnotesize\texttt{drove\_heavy} & \scriptsize yes/no & \footnotesize ... a truck or bus?\par\vspace{1pt}{\footnotesize\color{hdr}{\hifont ...ट्रक या बस?}} & \scriptsize 0 No | 1 Yes\par\vspace{1pt}\textit{Asked only if: Ask only if tk\_drive is 1 or 2} & \scriptsize Level ladder without licences\\
\footnotesize 183 & \footnotesize\texttt{read\_at\_work} & \scriptsize yes/no & \footnotesize In your work, do you read anything (notes, rate lists, tickets, messages)?\par\vspace{1pt}{\footnotesize\color{hdr}{\hifont अपने काम में क्या आप कुछ पढ़ते हैं (पर्ची, रेट लिस्ट, टिकट, मैसेज)?}} & \scriptsize 0 No | 1 Yes & \scriptsize STEP module 6A\\
\footnotesize 184 & \footnotesize\texttt{calc\_at\_work} & \scriptsize yes/no & \footnotesize In your work, do you work out prices or costs?\par\vspace{1pt}{\footnotesize\color{hdr}{\hifont अपने काम में क्या आप दाम या ख़र्चा जोड़ते-घटाते हैं?}} & \scriptsize 0 No | 1 Yes & \scriptsize STEP module 6A\\
\bottomrule\end{longtable}\vspace{6pt}
\noindent{\large\bfseries\color{hdr} Module M --- Access to common resources (short block)}\quad{\small(9 questions)}\\[3pt]
{\small\itshape Read aloud before this module: Last, a few questions about the forest and common land around where you work — firewood, fodder, grazing and water. Some of this is not allowed in some places, and that is a perfectly good answer: we only want to know what is possible and what is not. Nothing you say goes to the Forest Department or to anyone else, and you can leave any of these out.}\\[5pt]
\begin{longtable}{@{}p{0.5cm}p{3.3cm}p{1.7cm}p{7.3cm}p{5.6cm}p{5.4cm}@{}}
\toprule \footnotesize\textbf{\#} & \footnotesize\textbf{Variable} & \footnotesize\textbf{Type} & \footnotesize\textbf{Question (English / Hindi)} & \footnotesize\textbf{Options {\&} skip} & \footnotesize\textbf{Source}\\ \midrule
\endhead
\footnotesize 185 & \footnotesize\texttt{cpr\_permitted} & \scriptsize select one & \footnotesize Where you work, is it allowed to take firewood, fodder or forest produce from the forest or common land? Ask what they understand the rule to be; do not correct them.\par\vspace{1pt}{\footnotesize\color{hdr}{\hifont आप जहाँ काम करते हैं, वहाँ जंगल या सार्वजनिक ज़मीन से लकड़ी, चारा या जंगल की चीज़ें लेना मना है या नहीं? वे जो नियम समझते हैं वही पूछें; उन्हें सुधारें नहीं।}} & \scriptsize 1 Yes, it is allowed | 0 No, it is not allowed | 2 Allowed only with a permit or a fee | 97 Don't know\par\vspace{1pt}\textit{Asked only if: May be left blank} & \scriptsize Project design (de jure access, 2026-10-05); Ostrom 1990 on the distinction between rule in form and rule in use\\
\footnotesize 186 & \footnotesize\texttt{cpr\_firewood} & \scriptsize select one & \footnotesize In this Yatra season or last Yatra season, did you or anyone staying with you here collect firewood or dead wood from common land or forest in the Yatra region or workplace region?\par\vspace{1pt}{\footnotesize\color{hdr}{\hifont इस यात्रा सीज़न या पिछले यात्रा सीज़न में, क्या आपने या आपके साथ यहाँ रहने वाले किसी ने यात्रा क्षेत्र या काम की जगह की साझा ज़मीन या जंगल से लकड़ी या सूखी टहनियाँ इकट्ठी कीं?}} & \scriptsize 0 No | 1 Yes | 2 No, it is not allowed here | 97 Don't know\par\vspace{1pt}\textit{Asked only if: May be left blank} & \scriptsize Krantz 2001, natural capital: resource stocks from which resource flows are derived (definitions list); item wording is ours\\
\footnotesize 187 & \footnotesize\texttt{cpr\_fodder} & \scriptsize select one & \footnotesize In this Yatra season or last Yatra season, did you or anyone staying with you here cut grass or fodder from common land in the Yatra region or workplace region?\par\vspace{1pt}{\footnotesize\color{hdr}{\hifont इस यात्रा सीज़न या पिछले यात्रा सीज़न में, क्या आपने या आपके साथ यहाँ रहने वाले किसी ने यात्रा क्षेत्र या काम की जगह की साझा ज़मीन से घास या चारा काटा?}} & \scriptsize 0 No | 1 Yes | 2 No, it is not allowed here | 97 Don't know\par\vspace{1pt}\textit{Asked only if: May be left blank} & \scriptsize Krantz 2001, natural capital: resource stocks from which resource flows are derived; item wording is ours\\
\footnotesize 188 & \footnotesize\texttt{cpr\_forest\_produce} & \scriptsize select one & \footnotesize In this Yatra season or last Yatra season, did you or anyone staying with you here collect herbs, mushrooms or wild fruit from the forest in the Yatra region or workplace region?\par\vspace{1pt}{\footnotesize\color{hdr}{\hifont इस यात्रा सीज़न या पिछले यात्रा सीज़न में, क्या आपने या आपके साथ यहाँ रहने वाले किसी ने यात्रा क्षेत्र या काम की जगह के जंगल से जड़ी-बूटी, मशरूम या जंगली फल इकट्ठे किए?}} & \scriptsize 0 No | 1 Yes | 2 No, it is not allowed here | 97 Don't know\par\vspace{1pt}\textit{Asked only if: May be left blank} & \scriptsize Krantz 2001, natural capital (resource flows from stocks); item wording is ours\\
\footnotesize 189 & \footnotesize\texttt{cpr\_grazing} & \scriptsize select one & \footnotesize In this Yatra season or last Yatra season, did any animal of yours or of anyone staying with you here graze on common pasture in the Yatra region or workplace region?\par\vspace{1pt}{\footnotesize\color{hdr}{\hifont इस यात्रा सीज़न या पिछले यात्रा सीज़न में, क्या आपका या आपके साथ यहाँ रहने वाले किसी का कोई जानवर यात्रा क्षेत्र या काम की जगह की साझा चरागाह में चरा?}} & \scriptsize 0 No | 1 Yes | 2 No, it is not allowed here | 97 Don't know\par\vspace{1pt}\textit{Asked only if: May be left blank} & \scriptsize Krantz 2001, natural capital: resource stocks from which resource flows are derived; item wording is ours\\
\footnotesize 190 & \footnotesize\texttt{cpr\_restricted} & \scriptsize select one & \footnotesize In this Yatra season or last Yatra season, has anyone you or someone staying with you here been stopped from using a common forest, pasture or path, or been asked to pay to use it, in the Yatra region or workplace region?\par\vspace{1pt}{\footnotesize\color{hdr}{\hifont इस यात्रा सीज़न या पिछले यात्रा सीज़न में, क्या आपको या आपके साथ यहाँ रहने वाले किसी को यात्रा क्षेत्र या काम की जगह में साझा जंगल, चरागाह या रास्ते के इस्तेमाल से रोका गया, या इसके लिए पैसे देने को कहा गया?}} & \scriptsize 0 No | 1 Yes | 97 Don't know\par\vspace{1pt}\textit{Asked only if: May be left blank} & \scriptsize Krantz 2001, access ('the opportunity in practice to use a resource') and claims ('demands and appeals'); item wording is ours\\
\footnotesize 191 & \footnotesize\texttt{cpr\_lost\_access} & \scriptsize select one & \footnotesize In this Yatra season or last Yatra season, has a common place you used to collect from or graze on been closed, because of a landslide, flood, road or bridge damage, or construction work, in the Yatra region or workplace region?\par\vspace{1pt}{\footnotesize\color{hdr}{\hifont इस यात्रा सीज़न या पिछले यात्रा सीज़न में, क्या कोई साझा जगह, जहाँ आप पहले लकड़ी, घास या चारा लेते थे, भूस्खलन, बाढ़, सड़क या पुल टूटने, या निर्माण कार्य की वजह से बंद हो गई?}} & \scriptsize 0 No | 1 Yes | 97 Don't know\par\vspace{1pt}\textit{Asked only if: May be left blank} & \scriptsize Krantz 2001, access (as above); the ropeway link is the project's own\\
\footnotesize 192 & \footnotesize\texttt{cpr\_water\_route} & \scriptsize select one & \footnotesize On your route or near your worksite, can you use a common water source (spring, tap or stream) without paying?\par\vspace{1pt}{\footnotesize\color{hdr}{\hifont रास्ते पर या काम की जगह के पास, क्या आप कोई साझा पानी का स्रोत (झरना, नल या नदी-नाला) बिना पैसे दिए इस्तेमाल कर सकते हैं?}} & \scriptsize 1 Yes, free to use | 2 A source exists but we must pay | 3 No source nearby | 97 Don't know\par\vspace{1pt}\textit{Asked only if: May be left blank} & \scriptsize Krantz 2001, natural capital (water) and access ('the opportunity in practice to use a resource'); item wording is ours\\
\footnotesize 193 & \footnotesize\texttt{cpr\_governance} & \scriptsize select one & \footnotesize In the Yatra region or workplace region, who decides who may use the common land and forest? Ask and code what they name; do not read the list.\par\vspace{1pt}{\footnotesize\color{hdr}{\hifont यात्रा क्षेत्र या काम की जगह की साझा ज़मीन और जंगल का फ़ैसला कौन करता है, कि कौन उसका इस्तेमाल कर सकता है? जो नाम लें वही लिख लें; सूची न पढ़ें।}} & \scriptsize 1 Village head or panchayat | 2 Forest department | 3 Local committee or community group | 4 Nobody decides; it is open to all | 97 Don't know\par\vspace{1pt}\textit{Asked only if: May be left blank} & \scriptsize NO SOURCE ON DISK YET: Krantz does not address governance of common land. To be sourced from Agrawal 2001 or Ostrom 1990 once read on disk\\
\bottomrule\end{longtable}\vspace{6pt}
\clearpage\noindent{\large\bfseries\color{hdr} Sources, and what each one contributed}\\[4pt]
\begin{longtable}{@{}p{9.5cm}p{1.2cm}p{12.5cm}@{}}
\toprule \footnotesize\textbf{Source} & \footnotesize\textbf{Items} & \footnotesize\textbf{Variables}\\ \midrule
\endhead
\scriptsize Standard household-survey item & \scriptsize 8 & \scriptsize\texttt{female, hoh\_relation, marital\_status, knows\_years\_schooling, migration\_referral\_other, work\_equipment\_detail, n\_life\_insured, has\_crop\_insurance}\\
\scriptsize FAO Food Insecurity Experience Scale (FIES), SDG indicator 2.1.2; behavioural items only & \scriptsize 5 & \scriptsize\texttt{fies\_skipped\_meal, fies\_ate\_less, fies\_ran\_out, fies\_hungry, fies\_whole\_day}\\
\scriptsize Chaudhuri et al. 2002 asset covariate & \scriptsize 4 & \scriptsize\texttt{owns\_shop\_stall, land\_cultivable\_acres, owns\_work\_vehicle, owns\_work\_equipment}\\
\scriptsize Project design (amount only; frequency dropped for length). A previous citation to NSS 64th Round practice was withdrawn 2026-09-28: no NSS schedule is held by this project and it could not be checked & \scriptsize 4 & \scriptsize\texttt{remit\_out\_yatra\_pm, remit\_out\_offseason\_pm, remit\_in\_yatra\_pm, remit\_in\_offseason\_pm}\\
\scriptsize Apablaza et al. 2026, Appendix 2 employment table, work intensity (1/16) & \scriptsize 3 & \scriptsize\texttt{intens\_speed, intens\_deadline, intens\_time}\\
\scriptsize Apablaza et al. 2026, Appendix 2 employment table, posture risk (1/16) & \scriptsize 3 & \scriptsize\texttt{posture\_position, posture\_loads, posture\_repetitive}\\
\scriptsize NITI Aayog National MPI (years of schooling); Chaudhuri et al. 2002 covariate & \scriptsize 2 & \scriptsize\texttt{years\_schooling, education\_level\_cat}\\
\scriptsize NITI Aayog National MPI school-attendance indicator & \scriptsize 2 & \scriptsize\texttt{n\_children\_6\_14, n\_children\_in\_school}\\
\scriptsize Apablaza et al. 2026 Q13 (decomposed into hours/day x days/week rather than asked as one weekly figure, for ease of recall) & \scriptsize 2 & \scriptsize\texttt{hours\_day\_yatra, days\_week\_yatra}\\
\scriptsize Project design (replaces the twelve-cell status calendar, 2026-10-05) & \scriptsize 2 & \scriptsize\texttt{months\_worked\_yatra, months\_no\_paid\_work}\\
\scriptsize HCES 2022-23, MoSPI Appendix A, Sections 5.1-5.3 (30-day); IHDS-II Income and Social Capital Questionnaire, Q14 14.1-14.6 (30-day, same items) & \scriptsize 2 & \scriptsize\texttt{cons\_staples\_yatra\_pm, cons\_staples\_offseason\_pm}\\
\scriptsize HCES 2022-23, MoSPI Appendix A, Sections 6.1-6.8 (7-day actual recall in HCES; asked here as a usual monthly figure instead, for the reason above) & \scriptsize 2 & \scriptsize\texttt{cons\_perishables\_yatra\_pm, cons\_perishables\_offseason\_pm}\\
\scriptsize Calvo and Dercon 2007; Eze and Iheonu 2025 (non-purchased food, aggregate approach) & \scriptsize 2 & \scriptsize\texttt{cons\_food\_own\_yatra\_pm, cons\_food\_own\_offseason\_pm}\\
\scriptsize HCES 2022-23, MoSPI Appendix A, Section 7.1 'served processed food' (7-day in HCES; usual monthly here) & \scriptsize 2 & \scriptsize\texttt{cons\_food\_out\_yatra\_pm, cons\_food\_out\_offseason\_pm}\\
\scriptsize HCES 2022-23, MoSPI Appendix A, Section 8.1 (30-day). IHDS-II Income and Social Capital Questionnaire, Q14 14.21-14.22. Nigeria GHS-Panel Wave 3, Household Questionnaire, Sections 10B/11 (30-day tier) & \scriptsize 2 & \scriptsize\texttt{cons\_fuel\_yatra\_pm, cons\_fuel\_offseason\_pm}\\
\scriptsize HCES 2022-23, MoSPI Appendix A, Sections 9.1-9.2 (30-day). IHDS-II Income and Social Capital Questionnaire, Q14 14.25-14.27. Nigeria GHS-Panel Wave 3, Household Questionnaire, Sections 10B/11 (30-day tier) & \scriptsize 2 & \scriptsize\texttt{cons\_routine\_misc\_yatra\_pm, cons\_routine\_misc\_offseason\_pm}\\
\scriptsize HCES 2022-23, MoSPI Appendix A, Sections 11.1-11.2 (30-day). IHDS-II Income and Social Capital Questionnaire, Q14 14.24/14.28. Nigeria GHS-Panel Wave 3, Household Questionnaire, Sections 10B/11 (30-day tier) & \scriptsize 2 & \scriptsize\texttt{cons\_transport\_comm\_yatra\_pm, cons\_transport\_comm\_offseason\_pm}\\
\scriptsize HCES 2022-23, MoSPI Appendix A, Section 11.4 (30-day). IHDS-II Income and Social Capital Questionnaire, Q14 14.30. Nigeria GHS-Panel Wave 3, Household Questionnaire, Sections 10B/11 (30-day tier) & \scriptsize 2 & \scriptsize\texttt{cons\_rent\_yatra\_pm, cons\_rent\_offseason\_pm}\\
\scriptsize HCES 2022-23, MoSPI Appendix A, Section 10.3 (30-day). IHDS-II Income and Social Capital Questionnaire, Q14 14.33 (30-day, out-patient) & \scriptsize 2 & \scriptsize\texttt{cons\_med\_nonhosp\_yatra\_pm, cons\_med\_nonhosp\_offseason\_pm}\\
\scriptsize HCES 2022-23, MoSPI Appendix A, Section 7.2 (7-day in HCES; usual monthly here, as for every other seasonal item) & \scriptsize 2 & \scriptsize\texttt{cons\_packaged\_food\_yatra\_pm, cons\_packaged\_food\_offseason\_pm}\\
\scriptsize HCES 2022-23, MoSPI Appendix A, Sections 12.1-12.3 (7-day in HCES; usual monthly here, and combined into one figure) & \scriptsize 2 & \scriptsize\texttt{cons\_pan\_tobacco\_yatra\_pm, cons\_pan\_tobacco\_offseason\_pm}\\
\scriptsize NITI Aayog National MPI housing indicator & \scriptsize 2 & \scriptsize\texttt{floor\_material, roof\_material}\\
\scriptsize NITI Aayog National MPI drinking-water indicator (30-minute round-trip rule) & \scriptsize 2 & \scriptsize\texttt{water\_on\_premises, water\_fetch\_minutes}\\
\scriptsize Project design & \scriptsize 2 & \scriptsize\texttt{morbidity\_cost\_15d, ropeway\_stance}\\
\scriptsize Ligon and Schechter 2003; Dercon and Krishnan 2000 (VER needs shock magnitude, not only incidence) & \scriptsize 2 & \scriptsize\texttt{shock\_work\_lost\_weeks, shock\_money\_spent}\\
\scriptsize Apablaza et al. 2026, Appendix 2 Q20 (split into injury and death) & \scriptsize 2 & \scriptsize\texttt{workplace\_injury\_12m, workplace\_death\_12m}\\
\scriptsize Apablaza et al. 2026, Appendix 2 employment table, physical risk (1/16) & \scriptsize 2 & \scriptsize\texttt{phys\_noise, phys\_temperature}\\
\scriptsize Level ladder without licences & \scriptsize 2 & \scriptsize\texttt{drove\_car, drove\_heavy}\\
\scriptsize STEP module 6A & \scriptsize 2 & \scriptsize\texttt{read\_at\_work, calc\_at\_work}\\
\scriptsize Krantz 2001, natural capital: resource stocks from which resource flows are derived; item wording is ours & \scriptsize 2 & \scriptsize\texttt{cpr\_fodder, cpr\_grazing}\\
\scriptsize Standard; sensitivity covariate in the VEP model (Azeem et al. 2016 structure) & \scriptsize 1 & \scriptsize\texttt{age}\\
\scriptsize Project design (language proxy for regional/socioeconomic origin; avoids a sensitive caste question and covers out-of-country migrants) & \scriptsize 1 & \scriptsize\texttt{native\_language}\\
\scriptsize Project design (the language item stands in for caste; an uncoded 'other' is a real loss) & \scriptsize 1 & \scriptsize\texttt{native\_language\_other}\\
\scriptsize Project design (distance bucket); NOT an NSS classification — the earlier citation to one was withdrawn 2026-09-28, no NSS schedule is held by this project & \scriptsize 1 & \scriptsize\texttt{origin}\\
\scriptsize Standard state/UT classification; needed to apply the right poverty line per respondent & \scriptsize 1 & \scriptsize\texttt{home\_state}\\
\scriptsize Rural and urban poverty lines (Sethu et al. 2024) & \scriptsize 1 & \scriptsize\texttt{home\_rural\_urban}\\
\scriptsize STEP module 2; Chaudhuri et al. 2002 adaptive-capacity covariate & \scriptsize 1 & \scriptsize\texttt{training\_received}\\
\scriptsize Chaudhuri et al. 2002; used to get per-capita consumption & \scriptsize 1 & \scriptsize\texttt{hhsize}\\
\scriptsize NITI Aayog National MPI, Years of Schooling (weight 1/6) & \scriptsize 1 & \scriptsize\texttt{any\_member\_6yr\_schooling}\\
\scriptsize Adult-equivalent scales (Claro et al. 2010; Awuni et al. 2023); NITI Aayog National MPI nutrition age range & \scriptsize 1 & \scriptsize\texttt{n\_children\_u6}\\
\scriptsize Dependency measure used in VEP studies (Azeem et al. 2016) & \scriptsize 1 & \scriptsize\texttt{n\_earners}\\
\scriptsize Apablaza et al. 2026, Appendix 2 Q4 (adapted to a single-respondent report: 'who is the main contributor' becomes 'are you the main contributor') & \scriptsize 1 & \scriptsize\texttt{main\_income\_earner}\\
\scriptsize Project quota design, mapped to NCO-2015 families & \scriptsize 1 & \scriptsize\texttt{occupation}\\
\scriptsize Project design (coarse quota group + verbatim detail, coded to NCO-2015 in the office) & \scriptsize 1 & \scriptsize\texttt{occupation\_detail}\\
\scriptsize PLFS / NSS employment status classification [source not held by this project; citation unverified as of 2026-09-28] & \scriptsize 1 & \scriptsize\texttt{employment\_type}\\
\scriptsize Apablaza et al. 2026, Appendix 2 Q8 (shortened to wage workers); categories ours & \scriptsize 1 & \scriptsize\texttt{wage\_employer}\\
\scriptsize Project design (multiple job-holding, common in the Yatra economy) & \scriptsize 1 & \scriptsize\texttt{other\_activity\_types}\\
\scriptsize Project design (multiple job-holding) & \scriptsize 1 & \scriptsize\texttt{other\_activity\_income\_pm}\\
\scriptsize Apablaza et al. 2026 Appendix 2 Q7 (concept: job stability within the engagement); asked as a count per this project's no-scale rule & \scriptsize 1 & \scriptsize\texttt{employers\_in\_season}\\
\scriptsize Apablaza et al. 2026, Appendix 2 Q10; Table 3 (stability) & \scriptsize 1 & \scriptsize\texttt{years\_current\_job}\\
\scriptsize Occupational mobility check; job-history item & \scriptsize 1 & \scriptsize\texttt{prev\_occ\_change}\\
\scriptsize Gathmann and Schonberg 2010 (observed moves are systematically shorter than random moves — the validation test for a task-distance measure); coded to NCO-2015 in the office & \scriptsize 1 & \scriptsize\texttt{prev\_occ}\\
\scriptsize Project design (stated destination, for the structural-displacement analysis); coded to NCO-2015 in the office & \scriptsize 1 & \scriptsize\texttt{target\_occ}\\
\scriptsize Job-history item & \scriptsize 1 & \scriptsize\texttt{prev\_occ\_reason}\\
\scriptsize Project design (replaces the modal non-Yatra calendar cell, 2026-10-05) & \scriptsize 1 & \scriptsize\texttt{offseason\_activity}\\
\scriptsize Apablaza et al. 2026 Q15 (annual-total fallback), adapted to a seasonal workforce and now the only earnings route & \scriptsize 1 & \scriptsize\texttt{income\_annual\_total}\\
\scriptsize Project design (the Yatra / non-Yatra split needed to reweight an annual total onto the season) & \scriptsize 1 & \scriptsize\texttt{pct\_income\_yatra}\\
\scriptsize Apablaza et al. 2026, Appendix 2 Q15; bands to check & \scriptsize 1 & \scriptsize\texttt{earnings\_range}\\
\scriptsize Apablaza et al. 2026 Q13, asked a second time for the off-season (decomposed as hours/day x days/week, as in the Yatra-season pair) & \scriptsize 1 & \scriptsize\texttt{hours\_day\_offseason}\\
\scriptsize Apablaza et al. 2026 Q13, asked a second time for the off-season & \scriptsize 1 & \scriptsize\texttt{days\_week\_offseason}\\
\scriptsize Project design; categories from the pilot's own residency\_pattern distribution (n=46) & \scriptsize 1 & \scriptsize\texttt{resp\_returns\_at\_closure}\\
\scriptsize Project design (split-household measurement, underpinning Module E) & \scriptsize 1 & \scriptsize\texttt{hh\_at\_home\_place}\\
\scriptsize Project design (mirrors spend\_differs\_by\_season, Module E) & \scriptsize 1 & \scriptsize\texttt{remit\_differs\_by\_season}\\
\scriptsize Project design (remittance channel as a financial-inclusion and cost measure) & \scriptsize 1 & \scriptsize\texttt{remit\_mode}\\
\scriptsize Project design (enterprise concentration of household labour, 2026-10-05) & \scriptsize 1 & \scriptsize\texttt{n\_hh\_same\_business}\\
\scriptsize Project design (duration of the seasonal relationship) & \scriptsize 1 & \scriptsize\texttt{years\_coming\_here}\\
\scriptsize Project design (push/pull as a vulnerability covariate) & \scriptsize 1 & \scriptsize\texttt{came\_here\_reason}\\
\scriptsize Project design (catch-all codes lose the case they catch) & \scriptsize 1 & \scriptsize\texttt{came\_here\_reason\_other}\\
\scriptsize Project design (off-season labour migration) & \scriptsize 1 & \scriptsize\texttt{worked\_away\_in\_closure}\\
\scriptsize Project design; coded to NCO-2015 in the office & \scriptsize 1 & \scriptsize\texttt{closure\_work\_detail}\\
\scriptsize Project design; horizoned intention item after VASyR 2025 (move\_accom\_yesno, asked over a stated 6-month horizon) & \scriptsize 1 & \scriptsize\texttt{would\_move\_for\_work}\\
\scriptsize Project design (referral channel as a proxy for social capital and for placement debt) & \scriptsize 1 & \scriptsize\texttt{migration\_referral}\\
\scriptsize Project design (respondent-gated seasonal recall); the underlying two-season design follows Beegle et al. 2012 and Deaton and Grosh 2000 on usual-period recall & \scriptsize 1 & \scriptsize\texttt{spend\_differs\_by\_season}\\
\scriptsize HCES 2022-23, MoSPI Appendix A, Sections 13.1-13.2 (365-day). IHDS-II Income and Social Capital Questionnaire, Q14 14.38-14.39 (365-day). Nigeria GHS-Panel Wave 3, Household Questionnaire, Sections 10B/11 instead uses a 6-month recall for clothing (items 401-441); we follow the two India sources, since one of them (HCES) underlies our poverty line & \scriptsize 1 & \scriptsize\texttt{cons\_clothing\_12m}\\
\scriptsize HCES 2022-23, MoSPI Appendix A, Section 10.1 (365-day). IHDS-II Income and Social Capital Questionnaire, Q14 14.35-14.37 (365-day, same items). Nigeria GHS-Panel Wave 3, Household Questionnaire, Sections 10B/11 asks school fees inside its education/roster module, per child, not as a household total, so it is not directly comparable & \scriptsize 1 & \scriptsize\texttt{cons\_education\_12m}\\
\scriptsize HCES 2022-23, MoSPI Appendix A, Section 10.2 (365-day, hospitalisation). IHDS-II Income and Social Capital Questionnaire, Q14 14.34 (365-day, in-patient). Nigeria GHS-Panel Wave 3, Household Questionnaire, Sections 10B/11 instead uses one combined 6-month medical item (less detail); we follow the matching HCES/IHDS split & \scriptsize 1 & \scriptsize\texttt{cons\_medical\_hosp\_12m}\\
\scriptsize HCES 2022-23, MoSPI Appendix A, Section 14 (365-day, 10 sub-groups). IHDS-II Income and Social Capital Questionnaire, Q14 14.40-14.48 (365-day). Nigeria GHS-Panel Wave 3, Household Questionnaire, Sections 10B/11 (12-month tier: durables, appliances, building materials). All three agree on the 12-month recall; insurance premiums, which IHDS (14.50) and the Nigeria survey (items 510-513) count here, are excluded, following HCES and standard national-accounts practice (insurance is a financial item, not consumption) — premiums are covered instead by the insurance items in Module G & \scriptsize 1 & \scriptsize\texttt{cons\_durables\_12m}\\
\scriptsize NITI Aayog National MPI housing indicator (wall limb) & \scriptsize 1 & \scriptsize\texttt{wall\_material}\\
\scriptsize Lyons et al. 2023, Table 2 (area/settlement conditions), adapted to a labour-migrant setting & \scriptsize 1 & \scriptsize\texttt{accom\_type\_here}\\
\scriptsize NITI Aayog National MPI & \scriptsize 1 & \scriptsize\texttt{electricity}\\
\scriptsize NITI Aayog National MPI sanitation indicator (improved/unimproved and shared/not) & \scriptsize 1 & \scriptsize\texttt{toilet\_type}\\
\scriptsize NITI Aayog National MPI drinking-water indicator (JMP improved/unimproved taxonomy) & \scriptsize 1 & \scriptsize\texttt{drinking\_water}\\
\scriptsize NITI Aayog National MPI cooking-fuel indicator (its own list of dirty fuels) & \scriptsize 1 & \scriptsize\texttt{cooking\_fuel}\\
\scriptsize Project design (local land units, Uttarakhand) & \scriptsize 1 & \scriptsize\texttt{land\_unit}\\
\scriptsize NITI Aayog National MPI, Assets (weight 1/21) & \scriptsize 1 & \scriptsize\texttt{assets\_owned}\\
\scriptsize Chaudhuri et al. 2002 asset covariate; pony/mule class is this project's exposure variable & \scriptsize 1 & \scriptsize\texttt{livestock\_owned}\\
\scriptsize NITI Aayog National MPI bank-account indicator; VEP adaptive-capacity covariate & \scriptsize 1 & \scriptsize\texttt{has\_bank\_account}\\
\scriptsize VEP adaptive-capacity covariate & \scriptsize 1 & \scriptsize\texttt{took\_loan\_12m}\\
\scriptsize VEP adaptive-capacity covariate (lender type) & \scriptsize 1 & \scriptsize\texttt{credit\_source}\\
\scriptsize Project design (investment against distress borrowing) & \scriptsize 1 & \scriptsize\texttt{loan\_purpose}\\
\scriptsize Project design (loan size, for the debt-burden measures) & \scriptsize 1 & \scriptsize\texttt{loan\_amount\_borrowed}\\
\scriptsize Islam and Chowdhury 2025 (financial distress and household vulnerability to poverty) & \scriptsize 1 & \scriptsize\texttt{loan\_amount}\\
\scriptsize Project design (informal lending is often interest-free) & \scriptsize 1 & \scriptsize\texttt{pays\_interest}\\
\scriptsize Project design (informal-credit pricing as locally quoted) & \scriptsize 1 & \scriptsize\texttt{loan\_interest\_per100\_pm}\\
\scriptsize Carter and Zimmerman 2000; Zimmerman and Carter 2003 (asset-based coping) & \scriptsize 1 & \scriptsize\texttt{loan\_collateral}\\
\scriptsize NITI Aayog National MPI health-insurance indicator; Lyons et al. 2023 & \scriptsize 1 & \scriptsize\texttt{n\_health\_insured}\\
\scriptsize Project design; scheme names as administered in Uttarakhand & \scriptsize 1 & \scriptsize\texttt{govt\_any\_benefit}\\
\scriptsize Project design; coded in the office against the Uttarakhand scheme list & \scriptsize 1 & \scriptsize\texttt{govt\_schemes\_detail}\\
\scriptsize VEP adaptive-capacity covariate; STEP digital items & \scriptsize 1 & \scriptsize\texttt{smartphone\_owned}\\
\scriptsize Project design (intra-household financial capability) & \scriptsize 1 & \scriptsize\texttt{n\_can\_transact\_online}\\
\scriptsize NSS health module (15-day recall) [source not held by this project; citation unverified as of 2026-09-28] & \scriptsize 1 & \scriptsize\texttt{morbidity\_15d}\\
\scriptsize Project design (a reported illness with no cost/coping follow-up said nothing about its burden) & \scriptsize 1 & \scriptsize\texttt{morbidity\_coping\_15d}\\
\scriptsize Washington Group on Disability Statistics, Short Set item 3 (mobility), verbatim; Lyons et al. 2023 health dimension & \scriptsize 1 & \scriptsize\texttt{func\_limitation}\\
\scriptsize NSS health module (365-day recall) [source not held by this project; citation unverified as of 2026-09-28] & \scriptsize 1 & \scriptsize\texttt{hospitalization\_365d}\\
\scriptsize VASyR 2025 barriers\_health\_case\_access\_phc\_m (primary-care access barriers); Lyons et al. 2023 Table 2 indicator 2 'healthcare access' & \scriptsize 1 & \scriptsize\texttt{health\_access\_barrier\_3m}\\
\scriptsize NITI Aayog National MPI, Child and Adolescent Mortality (weight 1/12) & \scriptsize 1 & \scriptsize\texttt{child\_death\_5y}\\
\scriptsize NITI Aayog National MPI, Maternal Health (weight 1/12) & \scriptsize 1 & \scriptsize\texttt{birth\_last\_5y}\\
\scriptsize NITI Aayog National MPI, Maternal Health (antenatal care limb) & \scriptsize 1 & \scriptsize\texttt{anc\_4\_visits}\\
\scriptsize NITI Aayog National MPI, Maternal Health (assisted delivery limb); skilled-provider definition per NFHS & \scriptsize 1 & \scriptsize\texttt{skilled\_birth\_attendant}\\
\scriptsize VEP exposure covariate (Azeem et al. 2016); shock inventory following Gunther and Harttgen 2009, via Fujii 2016 section 4 & \scriptsize 1 & \scriptsize\texttt{distress\_event\_last365d}\\
\scriptsize Dercon and Krishnan 2000 (seasonal timing of shocks); project design & \scriptsize 1 & \scriptsize\texttt{shock\_month}\\
\scriptsize Coping strategies in VEP studies; Carter and Zimmerman 2000 on asset versus consumption smoothing & \scriptsize 1 & \scriptsize\texttt{shock\_coping}\\
\scriptsize Pelz et al. 2024, Energy Policy 186:113973 (concept: perceived effect on livelihoods; wording ours) & \scriptsize 1 & \scriptsize\texttt{ropeway\_expect\_work}\\
\scriptsize Pelz et al. 2024, Energy Policy 186:113973 (concept: local vs outside benefit; wording ours) & \scriptsize 1 & \scriptsize\texttt{ropeway\_expect\_jobs}\\
\scriptsize Pelz et al. 2024, Energy Policy 186:113973 (Fig. 4B, trust in the company; wording ours) & \scriptsize 1 & \scriptsize\texttt{ropeway\_trust}\\
\scriptsize Project design (see tasks\_module W4: exposure must be asked directly) & \scriptsize 1 & \scriptsize\texttt{trek\_dependent}\\
\scriptsize Apablaza et al. 2026, Appendix 2 Q11 & \scriptsize 1 & \scriptsize\texttt{contract\_status}\\
\scriptsize Apablaza et al. 2026, Appendix 2 Q12 (widened to local registrations) & \scriptsize 1 & \scriptsize\texttt{workplace\_registered}\\
\scriptsize Apablaza et al. 2026, Appendix 2 Q16 & \scriptsize 1 & \scriptsize\texttt{pension\_contrib}\\
\scriptsize Apablaza et al. 2026, Appendix 2 Q17 & \scriptsize 1 & \scriptsize\texttt{work\_health\_ins}\\
\scriptsize Apablaza et al. 2026, Appendix 2 Q18 & \scriptsize 1 & \scriptsize\texttt{leave\_rights}\\
\scriptsize Apablaza et al. 2026, Appendix 2 Q19 & \scriptsize 1 & \scriptsize\texttt{injured\_ever}\\
\scriptsize Apablaza et al. 2026, Appendix 2 Q21 & \scriptsize 1 & \scriptsize\texttt{wants\_more\_work}\\
\scriptsize Apablaza et al. 2026, Appendix 2 Q22 & \scriptsize 1 & \scriptsize\texttt{more\_hours\_day}\\
\scriptsize Apablaza et al. 2026, Appendix 2 Q23 (duration in weeks, tied to the calendar instead of a single point in time) & \scriptsize 1 & \scriptsize\texttt{months\_looked\_for\_work}\\
\scriptsize Apablaza et al. 2026, Appendix 2 Q24 & \scriptsize 1 & \scriptsize\texttt{first\_job\_ever}\\
\scriptsize Apablaza et al. 2026, Appendix 2 employment table, occupational status (1/8) & \scriptsize 1 & \scriptsize\texttt{higher\_ed}\\
\scriptsize Apablaza et al. 2026, Appendix 2 employment table, employment security (1/8) & \scriptsize 1 & \scriptsize\texttt{social\_security\_any}\\
\scriptsize Gathmann and Schonberg 2010; NCO 9333.0100 loader & \scriptsize 1 & \scriptsize\texttt{tk\_load}\\
\scriptsize Spitz-Oener 2006 / STEP; NCO 8322.0100 & \scriptsize 1 & \scriptsize\texttt{tk\_drive}\\
\scriptsize Gathmann and Schonberg 2010; NCO 8343.1700 ropeway operator & \scriptsize 1 & \scriptsize\texttt{tk\_engine}\\
\scriptsize Spitz-Oener 2006; NCO 7411.0301 wireman & \scriptsize 1 & \scriptsize\texttt{tk\_electric}\\
\scriptsize Gathmann and Schonberg 2010; NCO 8343.1700 & \scriptsize 1 & \scriptsize\texttt{tk\_safety}\\
\scriptsize Gathmann and Schonberg 2010; NCO 5223.0200 & \scriptsize 1 & \scriptsize\texttt{tk\_sell}\\
\scriptsize Autor and Handel 2013; NCO 5131.0401 & \scriptsize 1 & \scriptsize\texttt{tk\_cash}\\
\scriptsize NCO 5120.0300 cook & \scriptsize 1 & \scriptsize\texttt{tk\_cook}\\
\scriptsize NCO 5131.0401 waiter; Pal et al. 2026 & \scriptsize 1 & \scriptsize\texttt{tk\_serve}\\
\scriptsize NCO 5151.0202 room attendant; Autor-Levy-Murnane 2003 (janitorial) & \scriptsize 1 & \scriptsize\texttt{tk\_clean}\\
\scriptsize Gathmann and Schonberg 2010; NCO 5113.0200 & \scriptsize 1 & \scriptsize\texttt{tk\_guide}\\
\scriptsize NCO 8343.1700 ropeway operator & \scriptsize 1 & \scriptsize\texttt{tk\_coord}\\
\scriptsize Level ladder without licences: people drive without holding a licence & \scriptsize 1 & \scriptsize\texttt{drove\_twowheeler}\\
\scriptsize Project design (de jure access, 2026-10-05); Ostrom 1990 on the distinction between rule in form and rule in use & \scriptsize 1 & \scriptsize\texttt{cpr\_permitted}\\
\scriptsize Krantz 2001, natural capital: resource stocks from which resource flows are derived (definitions list); item wording is ours & \scriptsize 1 & \scriptsize\texttt{cpr\_firewood}\\
\scriptsize Krantz 2001, natural capital (resource flows from stocks); item wording is ours & \scriptsize 1 & \scriptsize\texttt{cpr\_forest\_produce}\\
\scriptsize Krantz 2001, access ('the opportunity in practice to use a resource') and claims ('demands and appeals'); item wording is ours & \scriptsize 1 & \scriptsize\texttt{cpr\_restricted}\\
\scriptsize Krantz 2001, access (as above); the ropeway link is the project's own & \scriptsize 1 & \scriptsize\texttt{cpr\_lost\_access}\\
\scriptsize Krantz 2001, natural capital (water) and access ('the opportunity in practice to use a resource'); item wording is ours & \scriptsize 1 & \scriptsize\texttt{cpr\_water\_route}\\
\scriptsize NO SOURCE ON DISK YET: Krantz does not address governance of common land. To be sourced from Agrawal 2001 or Ostrom 1990 once read on disk & \scriptsize 1 & \scriptsize\texttt{cpr\_governance}\\
\bottomrule\end{longtable}
\end{document}
